% Draft manuscript source for the Q3LOCK phase-coexistence paper.
% PDF generation is intentionally deferred until content review, final
% organization, hash capture, and signed external review are complete.
\documentclass[11pt]{article}
\usepackage[a4paper,margin=1in]{geometry}
\usepackage{amsmath,amssymb,amsthm,mathtools}
\usepackage{booktabs}
\usepackage{array}
\usepackage{enumitem}
\usepackage[T1]{fontenc}
\usepackage{lmodern}
\usepackage{hyperref}

\hypersetup{colorlinks=true,linkcolor=blue,citecolor=blue,urlcolor=blue}
\numberwithin{equation}{section}
\newtheorem{theorem}{Theorem}[section]
\newtheorem{proposition}[theorem]{Proposition}
\newtheorem{lemma}[theorem]{Lemma}
\newtheorem{corollary}[theorem]{Corollary}
\theoremstyle{definition}
\newtheorem{definition}[theorem]{Definition}
\newtheorem{remark}[theorem]{Remark}
\newcommand{\R}{\mathbb R}
\newcommand{\Z}{\mathbb Z}
\newcommand{\T}{\mathbb T}
\newcommand{\E}{\mathbb E}
\newcommand{\dd}{\,\mathrm d}
\newcommand{\norm}[1]{\left\lVert#1\right\rVert}
\newcommand{\ip}[2]{\left\langle#1,#2\right\rangle}
\newcommand{\Var}{\operatorname{Var}}
\newcommand{\Tr}{\operatorname{Tr}}
\newcommand{\cG}{\mathcal G}
\newcommand{\eps}{\varepsilon}

\title{Phase Coexistence in an Eight-Component Q3-Locked Quantum Anharmonic Crystal}
\author{Jusang Lee}
\date{Draft v0.1.7 -- internal content review -- 2026-09-07}

\begin{document}
\maketitle

\begin{abstract}
We formulate a fixed-spacing three-dimensional quantum anharmonic crystal
with eight real components per lattice site, a ferromagnetic spatial coupling,
and a non-radial quartic locking interaction on the internal cube graph
$Q_3$.  The source is coupled to the normalized collective direction
$u=8^{-1/2}(1,\ldots,1)$.  This manuscript assembles a conditional theorem
package, with every imported hypothesis displayed, for the following chain:
finite-volume heat traces and Euclidean loop measures; open-box and periodic
thermodynamic pressure; compact tempered DLR states; continuous-loop FKG
association; spatial reflection positivity and a three-dimensional infrared
bound; a volume-uniform collective Duhamel lower bound; and, in an explicit
low-temperature sufficient regime, a strict source cusp and two distinct
parity-related zero-source DLR states.  The proof uses the general-vector
existence results of Kozitsky--Pasurek, the finite-dimensional infrared
inequality of Froehlich--Simon--Spencer through an explicit hypothesis
crosswalk.  The needed Griffiths endpoint specialization and finite spectral
Falk--Bruch inequality are proved directly.  No scalar or
radial phase theorem is imported.  The result is currently an internal
manuscript-content draft at tier T0 (R-497). The spatial DLR, infrared,
operator and final composition arguments have been expanded, but signed mathematical review,
independent literature review, a clean-snapshot replay, and final content
freeze remain open.  In particular, this paper makes no claim about a ground
state gap, real-time KMS dynamics, a continuum limit, a physical vacuum,
cosmology, or closure of any TECT sector.  PDF generation and submission are
deliberately deferred until those gates are complete.
\end{abstract}

\section{Scope, status, and contributions}

This paper concerns one explicitly defined mathematical model.  It is not a
claim that the model is the microscopic law of nature.  The research record
for the assembled content is EXP-000780--EXP-000782 and R-497.  R-497 is
registered as \texttt{T0}, \texttt{claim\_bearing=false}, and
\texttt{INTERNAL\_REVIEW\_ONLY}.  The statements below are therefore written
as a conditional composition while the model-specific arguments and their
source applicability are subjected to line-by-line independent review.

The paper has seven proof blocks.
\begin{enumerate}[label=\textup{(B\arabic*)},leftmargin=2.3em]
\item finite-volume Hamiltonians, heat-trace finiteness, and the exact
      harmonic-reference/path-integral normalization;
\item open-box pressure, periodic pressure, and locally uniform source
      convergence;
\item tempered Euclidean DLR existence, compactness, and source-tangent
      limits;
\item continuous-loop FKG association and the collective covariance
      inequalities;
\item spatial reflection positivity, finite-mesh FSS Gaussian domination,
      and the nonzero-mode infrared sum in dimension three;
\item a collective translation, form-domain, and Falk--Bruch lower bound;
\item strict cusp and construction of a parity-related DLR pair under an
      explicit sufficient condition.
\end{enumerate}

The novelty boundary is intentionally narrow.  The standard ingredients are
the cited Gibbs/DLR, FKG, infrared, and Griffiths tools.  The model-specific
content is their typed composition for a non-radial eight-component $Q_3$
interaction, including the exact source normalization and the collective
lower bound.  No priority or world-first assertion is made.

The following conclusions are explicitly outside the paper:
\begin{itemize}[leftmargin=2.3em]
\item no all-parameter phase theorem, and no phase-absence theorem outside
      the sufficient regime below;
\item no assertion that the two states are extremal, pure, clustering,
      unique, or the complete DLR simplex;
\item no common infinite-volume real-time dynamics, algebraic KMS statement,
      ground-state phase, or spectral gap;
\item no spatial lattice-spacing removal, continuum renormalization, physical-vacuum
      ordering, cosmological interpretation, or Yang--Mills conclusion;
\item no closure of TECT sectors A, C, CP1, Pre-A, or the master theorem;
\item no submission, upload, release tag, or PDF before the gates in
      Section~\ref{sec:readiness} are discharged.
\end{itemize}

\section{The exact Q3LOCK model and conventions}
\label{sec:model}

Let $\Lambda_L=(\Z/L\Z)^3$, $V=L^3$, and take even $L\geq4$ in the
thermodynamic sequence.  Positive-direction nearest-neighbour bonds are
counted once for $L\geq4$; the finite-volume source record also specifies the
parallel-bond multiplicity convention at $L=2$.  A site variable is
$q_y=(q_{y,e})_{e\in\{0,1\}^3}\in\R^8$.  The internal graph $Q_3$ has an
edge between two binary triples that differ in one coordinate.  Write
\[
 u=8^{-1/2}(1,\ldots,1),\qquad Q_y=(u,q_y),\qquad S_y=|q_y|^2,
\]
and
\[
 D_y^{\mathrm{int}}=\sum_{\{e,f\}\in E(Q_3)}(q_{y,e}-q_{y,f})^2.
\]
The superscript distinguishes this internal graph quantity from the Duhamel
matrix introduced later.

For $m=\chi/\hbar^2>0$, $c,g,\lambda>0$, $r\in\R$, and real collective
source $h$, define
\begin{align}
 U_{L,h}(q)
 &=\sum_{y\in\Lambda_L}\left[
       \frac r2S_y+\frac g4\sum_e q_{y,e}^4
       +\frac\lambda4\sum_{\{e,f\}\in E(Q_3)}
          (q_{y,e}-q_{y,f})^2(q_{y,e}^2+q_{y,f}^2)
       -hQ_y\right]\nonumber\\
 &\qquad+\frac c2\sum_{\langle y,z\rangle}|q_y-q_z|^2. \label{eq:potential}
\end{align}
No radial or $O(8)$ symmetry is assumed.  The Hamiltonian on
$L^2(\R^{8V},\dd q)$ is
\begin{equation}
 H_{L}(h)=-\frac1{2m}\sum_{y,e}\partial_{q_{y,e}}^2+U_{L,h}(q).
 \label{eq:hamiltonian}
\end{equation}
The phase application uses $r<0$, but the finite-pressure statements allow
arbitrary real $r$.

The closed quadratic-form domain is
\begin{equation}
 \mathsf Q_L=H^1(\R^{8V})\cap
 L^2\left(\R^{8V},\sum_y|q_y|^4\dd q\right).
 \label{eq:form-domain}
\end{equation}
The quartic scalar term and the nonnegative $Q_3$ and spatial terms dominate
all negative quadratic and linear source terms after a constant shift.  The
form is consequently semibounded and has compact form embedding at fixed
volume.  These are finite-dimensional facts; they do not constitute a
continuum or infinite-volume operator theorem.

\begin{remark}[Source and units]
The parameter $h$ is an energy coefficient in $H_L(h)=H_L(0)-h\sum_yQ_y$.
It is not $\beta h$.  This convention is responsible for every factor of
$\beta$ in the pressure and Duhamel identities below.
\end{remark}

\section{Finite volume: trace, loops, and harmonic split}
\label{sec:finite}

Choose an auxiliary $a>0$, set $\omega_a=\sqrt{a/m}$, and write
\begin{equation}
 H_a=-\frac1{2m}\Delta+\frac a2\sum_y|q_y|^2,\qquad
 R_{a,h}=U_{L,h}-\frac a2\sum_y|q_y|^2.
 \label{eq:harmonic-split}
\end{equation}
For $|h|\leq h_0$, put $b_a=|r-a|/2$ and
\begin{equation}
 A=\frac g{128},\qquad
 C=\frac{16b_a^2}{g}+\frac34 h_0^{4/3}(32/g)^{1/3},\qquad
 R_{a,h}(q)\geq A\sum_y|q_y|^4-VC.
 \label{eq:residual-lower}
\end{equation}
Indeed $\sum_e q_e^4\geq|q|^4/8$, the locking and spatial differences are
nonnegative, and
$b_at^2\leq gt^4/64+16b_a^2/g$ and
$h_0t\leq gt^4/128+(3/4)h_0^{4/3}(32/g)^{1/3}$ for $t\geq0$.
Min--max comparison with the
harmonic oscillator therefore yields
\begin{equation}
 0<Z_L(h):=\Tr e^{-\beta H_L(h)}
 \leq e^{\beta VC}\,[2\sinh(\beta\omega_a/2)]^{-8V}<\infty.
 \label{eq:trace-bound}
\end{equation}
This is an eigenvalue comparison, not an appeal to operator monotonicity of
the exponential.

For an even time mesh $N\geq4$, $\eps=\beta/N$, and cyclic variables
$x_{y,k}$, the Gaussian reference action is
\begin{equation}
 S_{a,N}(x)=\frac12\sum_{y,k}\left[
 \frac m\eps|x_{y,k+1}-x_{y,k}|^2+a\eps|x_{y,k}|^2\right].
 \label{eq:mesh-action}
\end{equation}
Let $n_L=8V$, and define the absolute Gaussian prefactor and probability law by
\[
 c_N=(m/(2\pi\eps))^{n_LN/2},\qquad
 Z_{a,N}=c_N\int e^{-S_{a,N}}\dd x,\qquad
 \gamma_{a,N,L}(\dd x)=\frac{c_N}{Z_{a,N}}e^{-S_{a,N}(x)}\dd x.
\]
The factor $c_N$ is the product of free heat-kernel normalizations.  For one
coordinate, Gaussian integration cancels it against the precision factor
$(m/\eps)^{N/2}$ and leaves the inverse square root of the cyclic determinant
with diagonal $2+t$, off-diagonal $-1$, and $t=\eps^2a/m$.
The exact cyclic determinant is
\begin{equation}
 Z_{a,N}=\left[2\sinh\left(N\operatorname{arsinh}
       \frac{\beta\omega_a}{2N}\right)\right]^{-8V},
 \qquad Z_{a,N}\longrightarrow [2\sinh(\beta\omega_a/2)]^{-8V}.
 \label{eq:cyclic-det}
\end{equation}
To see this, put $z+z^{-1}=2+t$, $z>1$.  Each Fourier eigenvalue is
$(z-\zeta)(z-\zeta^{-1})/z$, with $\zeta^N=1$. Their product is
$z^N+z^{-N}-2=4\sinh^2(N\log z/2)$ and
$\log z=2\operatorname{arsinh}(\sqrt t/2)$, proving \eqref{eq:cyclic-det}.
The independent recurrence is
$T_0=2$, $T_1=2+t$, $T_j=(2+t)T_{j-1}-T_{j-2}$, with determinant $T_N-2$.

\subsection{Gaussian convergence in the finite-volume loop space}
\label{sec:gaussian-loop}

All limits in this subsection keep $L,\beta,m,a$ fixed.  Write
$\mathsf E_L=C_{\rm per}([0,\beta];\R^{n_L})$, with its sup norm, and let
$\mathcal I_N$ be periodic linear interpolation.  For one scalar coordinate,
unitary discrete Fourier inversion gives the grid covariance
\begin{align}
 G_N(j)&=\frac1\beta\sum_{n\in\mathcal J_N}
       \frac{e^{2\pi i n j/N}}{D_N(n)},\qquad
 \mathcal J_N=\{-N/2+1,\ldots,N/2\},\nonumber\\
 D_N(n)&=a+\frac{4m}{\eps^2}\sin^2(\pi n/N),\qquad
 D(n)=a+\frac{4\pi^2mn^2}{\beta^2},\nonumber\\
 G(t)&=\frac1\beta\sum_{n\in\Z}\frac{e^{2\pi int/\beta}}{D(n)}.
 \label{eq:gaussian-covariances}
\end{align}
The continuum series is absolutely summable. It is positive definite and
defines the finite-dimensional distributions of a centered periodic Gaussian.
The increment estimate below provides a continuous version.  Both diagonal
covariances are bounded by
\begin{equation}
 G_N(0),G(0)\leq K_a:=\frac1{\beta a}+\frac\beta{12m}.
 \label{eq:gaussian-diagonal-bound}
\end{equation}
For the discrete bound, keep the zero Fourier mode and remove $a$ only from
the remaining denominators.  The cycle Laplacian pseudoinverse has diagonal
$(N^2-1)/(12N)$: its effective resistance at separation $j$ is
$j(N-j)/N$, whose sum is twice $N$ times that diagonal.  Multiplication by
$\eps/m$ gives the bound.  In the continuum, sum the nonzero comparison
coefficients $\beta/(4\pi^2mn^2)$ to obtain the same upper bound.

For $0<|n|\leq N/2$, the inequalities
$\sin x\geq2x/\pi$ and $0\leq x^2-\sin^2x\leq x^4/3$ on
$[0,\pi/2]$ imply
\[
 0\leq D_N(n)^{-1}-D(n)^{-1}
 \leq\frac{\pi^2\beta^2}{48mN^2}.
\]
The omitted modes obey $\sum_{n\notin\mathcal J_N}n^{-2}\leq6/N$;
one may sum the two decreasing tails by integrals and retain the missing
negative Nyquist mode separately. Consequently
\begin{equation}
 \sup_j|G_N(j)-G(\eps j)|
 \leq B_N:=\frac\beta{mN}
            \left(\frac{\pi^2}{48}+\frac{3}{2\pi^2}\right).
 \label{eq:gaussian-grid-error}
\end{equation}

Let $d$ denote circular distance between two times. The continuum
massless comparison, summing the Fourier series for the periodic bridge
increment, gives
$\E|X(t)-X(s)|^2\leq d(\beta-d)/(m\beta)\leq d/m$.
Cauchy--Schwarz then gives
$|G(t)-G(s)|\leq\sqrt{K_a d/m}$.
The interpolated covariance is a convex combination of four grid
covariances, and the corresponding time differences are within $2\eps$
of $t-s$.  Therefore
\begin{equation}
 \sup_{s,t}|\operatorname{Cov}(\mathcal I_NX(t),\mathcal I_NX(s))-G(t-s)|
 \leq B_N+\sqrt{\frac{2K_a\beta}{mN}}\longrightarrow0.
 \label{eq:gaussian-interpolated-error}
\end{equation}
This proves convergence of every finite Gaussian evaluation vector.

Tightness requires an estimate at arbitrary times, not just grid vertices.
For a grid increment at shortest separation $j$, comparison on the
zero-sum subspace with the massless cycle gives
$\E|x_k-x_l|^2\leq\eps j/m$.
If $d\leq\eps$, the interpolated increment crosses at most two partial
edges; the $L^2$ triangle inequality bounds its variance by
$d^2/(m\eps)\leq d/m$. If $d>\eps$, insert neighboring vertices at
both ends. The two end increments have $L^2$ norms at most
$\sqrt{\eps/m}$, and the intervening grid increment has norm at most
$\sqrt{(d+2\eps)/m}$. It follows that
\begin{align}
 \E|\mathcal I_NX(t)-\mathcal I_NX(s)|^2&\leq14d/m,\nonumber\\
 \E\|\mathcal I_NX(t)-\mathcal I_NX(s)\|^4
 &\leq n_L(n_L+2)14^2d^2/m^2.
 \label{eq:gaussian-tightness}
\end{align}
Here $(2+\sqrt3)^2<14$, and the second line uses independence of the
$n_L$ coordinates and the exact fourth moment of a centered Gaussian.
Kolmogorov's tightness criterion, with fourth-moment exponent $2>1$ and
the tight one-time marginals from \eqref{eq:gaussian-diagonal-bound},
applies on the time interval. Periodicity is a closed condition. Tightness
and \eqref{eq:gaussian-interpolated-error} thus prove
\begin{equation}
 \nu_N:=(\mathcal I_N)_*\gamma_{a,N,L}
       \ \Longrightarrow\ \gamma_{a,L}\quad\hbox{on }\mathsf E_L.
 \label{eq:gaussian-weak-limit}
\end{equation}
Every subsequential limit has the same Gaussian finite-dimensional laws,
which determine a measure on this continuous-path space.  This proves
convergence of the whole sequence.  No spatial weighted topology or
uniformity as $L\to\infty$ has been used.

\subsection{Interacting weights, normalizers and source moments}
\label{sec:weighted-loop}

The full mesh action is independent of $a$ after cancellation:
\begin{equation}
 S_{a,N}(x)+\eps\sum_kR_{a,h}(x_k)
 =\frac m{2\eps}\sum_{y,k}|x_{y,k+1}-x_{y,k}|^2
   +\eps\sum_kU_{L,h}(x_k). \label{eq:split-cancellation}
\end{equation}
Define on all of $\mathsf E_L$
\[
 \mathcal R_{N,h}(\omega)=\eps\sum_kR_{a,h}(\omega(k\eps)),\qquad
 \mathcal R_h(\omega)=\int_0^\beta R_{a,h}(\omega(t))\dd t,
 \qquad w_{N,h}=e^{-\mathcal R_{N,h}},\quad w_h=e^{-\mathcal R_h}.
\]
These expressions agree with the mesh residual on polygonal paths.
Equation~\eqref{eq:residual-lower} bounds both weights above by
$M_L=e^{\beta VC}$.  Under either Gaussian reference, Wick averaging
gives the uniform upper bound $\E\mathcal R_{N,h},\E\mathcal R_h
\leq\beta V C_J^{\rm mesh}$, where
\[
 C_J^{\rm mesh}=4|r-a|K_a+24cK_a+(6g+24\lambda)K_a^2.
\]
Indeed a scalar fourth moment is $3s^2$, an internal edge polynomial
has mean $8s^2$, there are twelve internal edges and at most three spatial
bonds per site, and $s\leq K_a$. The source has zero Gaussian mean.
Jensen's inequality consequently gives
\begin{equation}
 \mathfrak z_{N,h}:=\int w_{N,h}\dd\nu_N\geq e^{-\beta VC_J^{\rm mesh}},
 \qquad \mathfrak z_h:=\int w_h\dd\gamma_{a,L}
                         \geq e^{-\beta VC_J^{\rm mesh}}.
 \label{eq:mesh-normalizer-lower}
\end{equation}
The constants are uniform in $N$ and $|h|\leq h_0$ at this fixed volume.

Let $\mathcal K\subset\mathsf E_L$ be compact. Its paths have bounded
range and a common modulus of continuity $\eta_{\mathcal K}$ tending to
zero. The finite-dimensional polynomial $R_{a,h}$ has a common Lipschitz
constant $L_{\mathcal K}$ on that range for $|h|\leq h_0$. Thus
\begin{equation}
 \sup_{\omega\in\mathcal K,\ |h|\leq h_0}
 |\mathcal R_{N,h}(\omega)-\mathcal R_h(\omega)|
 \leq\beta L_{\mathcal K}\eta_{\mathcal K}(\eps)\longrightarrow0.
 \label{eq:compact-riemann-residual}
\end{equation}
The kinetic action is already included in $\nu_N$; it is not evaluated as
a Riemann sum of a derivative of a typical Gaussian path.

For any bounded continuous $F$, choose a common compact set carrying all
but probability $\eta$ of every $\nu_N$, using
\eqref{eq:gaussian-weak-limit}. On this set the weights converge uniformly
by \eqref{eq:compact-riemann-residual}; on its complement,
$|F(w_{N,h}-w_h)|\leq2\|F\|_\infty M_L$.  Finally $Fw_h$ is bounded
continuous, so ordinary weak convergence applies to it. Taking $N\to\infty$
and then $\eta\downarrow0$ proves
\begin{equation}
 \int Fw_{N,h}\dd\nu_N\longrightarrow
       \int Fw_h\dd\gamma_{a,L},\qquad
 \mathfrak z_{N,h}\longrightarrow\mathfrak z_h>0,
 \qquad
 \mu_{N,L,h}:=\frac{w_{N,h}\nu_N}{\mathfrak z_{N,h}}
       \Longrightarrow\mu_{L,h}:=\frac{w_h\gamma_{a,L}}{\mathfrak z_h}.
 \label{eq:weighted-loop-limit}
\end{equation}
This is weak convergence on $\mathsf E_L$, not total-variation convergence.
The absolute mesh partition function is $Z_{N,L}(h)=Z_{a,N}\mathfrak z_{N,h}$,
so it converges to $Z_a\mathfrak z_h$.  The full finite-mesh action
\eqref{eq:split-cancellation} and its prefactor $c_N$ do not depend on $a$.
Repeating the limit for any other fixed positive split therefore proves
independence of both this absolute limit and the normalized loop law from
the chosen split.

\subsection{Identification with the heat trace}
\label{sec:fk-identification}

The precise external input is Simon \cite[Theorem 1.1, (1.1)--(1.3),
page 2]{Simon-FK}: the Brownian-bridge formula for a continuous potential
bounded below on a finite-dimensional Euclidean space. Its kinetic
operator is $-\Delta/2$. With $D=8|\Lambda_L|$ and $z=\sqrt m\,q$,
the unitary map $\psi(q)\mapsto m^{-D/4}\psi(z/\sqrt m)$ sends our
operator to $-\Delta_z/2+U_{L,h}(z/\sqrt m)$. The polynomial $U_{L,h}$ is
continuous and bounded below at every fixed $L,h$, by the retained quartic
bound. In the original coordinates the free bridge has covariance
$m^{-1}(\min(s,t)-st/\beta)$, and its kernel prefactor is
$(m/(2\pi\beta))^{D/2}$. No mass-one convention is silently imposed.
The stronger unbounded-semigroup Theorem 1.2 of that source is not used.

The harmonic-reference version follows from this free-bridge input by defining
\begin{equation}
 \dd\mathbb P^a_{q,q'}=
 \frac{\exp[-\int_0^\beta a|\omega(t)|^2/2\,\dd t]}
 {\E^0_{q,q'}\exp[-\int_0^\beta a|\omega(t)|^2/2\,\dd t]}
 \dd\mathbb P^0_{q,q'}.
 \label{eq:fk-harmonic-bridge-density}
\end{equation}
The denominator is strictly positive. Applying the free formula to the
harmonic potential alone identifies it with $K_a/K_0$. Applying the same
formula to $U_{L,h}=a|q|^2/2+R_{a,h}$ gives the following identity.
For the positive harmonic operator $H_a$, a continuous real multiplication
potential $R$ bounded below, and $H=H_a+R$ defined by its semibounded form,
its kernel is
\[
 K_H(\beta;q,q')=K_a(\beta;q,q')
       \E^a_{q,q'}\exp\left[-\int_0^\beta R(\omega(t))\dd t\right].
\]
Here $\E^a_{q,q'}$ denotes the normalized harmonic bridge. The bounded
potential formula extends to $R$ by truncating $R$ above and passing by
monotone form convergence; the bridge factors are dominated by
$e^{-\beta\inf R}$. The kernels are continuous by local endpoint dominated
convergence in the bridge formula. Their semigroup property and symmetry give
$\int K_H(\beta;q,q)\dd q=\|e^{-\beta H/2}\|_{\rm HS}^2
=\Tr e^{-\beta H}$ by Tonelli; finiteness follows also from the
integrable majorant $e^{-\beta\inf R}K_a(\beta;q,q)$.
Mixing the diagonal bridges with $K_a(\beta;q,q)\dd q/Z_a$ gives
the centered periodic Gaussian with covariance
$(-m\partial_t^2+a)^{-1}$ in \eqref{eq:gaussian-covariances}.
Applying this formula with $R=R_{a,h}$ gives
\begin{equation}
 Z_L(h)=Z_a\,\E_{\gamma_{a,L}}
 \exp\left[-\int_0^\beta R_{a,h}(\omega(\tau))\dd\tau\right],
 \label{eq:fk}
\end{equation}
The parameter and form hypotheses are verified in
\eqref{eq:form-domain} and \eqref{eq:residual-lower}.  The source-specific
normalized Feynman--Kac formulas being compared are KP~\cite{KP},
Section 2.3, equations (2.28)--(2.32). Their $Z_\Lambda$ in (2.32) is
the residual normalizer $\mathfrak z_h$; the absolute heat trace here includes
$Z_a$. This identifies the weak limit above with the operator loop law and
supplies the absolute normalization without inferring it from a normalized
correlation identity. Simon's free-bridge theorem is the external analytic
input; the mass scaling, harmonic reweighting and trace normalization above
are its finite-volume application, not additional infinite-volume imports.

\subsection{Unbounded sources and pressure derivatives}
\label{sec:source-moments}

For a fixed real site field $f$, define
\[
 Y_f(\omega)=\int_0^\beta\sum_yf_y(u,\omega_y(t))\dd t,\qquad
 Y_{N,f}(\omega)=\eps\sum_{y,k}f_y(u,\omega_y(k\eps)),\qquad
 \mathcal Q_N(\omega)=\eps\sum_{y,k}|\omega_y(k\eps)|^4.
\]
On polygonal paths $Y_f=Y_{N,f}$ exactly: the integral of each affine
segment is the average of its endpoints, and cyclic summation counts each
endpoint twice. Holder's inequality gives
$|Y_{N,f}|\leq B_f\mathcal Q_N^{1/4}$, where
$B_f=(\beta\sum_y|f_y|^{4/3})^{3/4}$.
For $b,\delta>0$, maximization over $z\geq0$ gives
\[
 bz-\delta z^4\leq C_Y(b,\delta)
 :=\frac{3b^{4/3}}{4(4\delta)^{1/3}}.
\]
Using \eqref{eq:residual-lower} and \eqref{eq:mesh-normalizer-lower},
for every $T>0$, $p>1$ and $0<\delta<A$, we obtain
\begin{equation}
 \sup_{N,\ |h|\leq h_0}\E_{\mu_{N,L,h}}e^{pT|Y_{N,f}|}
 \leq\exp\{\beta V(C+C_J^{\rm mesh})+C_Y(pTB_f,\delta)\}.
 \label{eq:mesh-source-ui}
\end{equation}
The analogous continuous estimate follows by integral Holder and the same
Young inequality. A slightly larger $T$ absorbs any fixed power of
$|Y_{N,f}|$. Truncation and \eqref{eq:weighted-loop-limit} therefore pass
the source exponential and its first two derivatives to the limit.

Point evaluations use a different bound. The density relative to the
Gaussian reference is at most $e^{\beta V(C+C_J^{\rm mesh})}$, and a
polygonal Gaussian evaluation has variance at most $K_a$. Hence
\begin{equation}
 \sup_{N,\ |h|\leq h_0,t,y,e}
 \E_{\mu_{N,L,h}}|\omega_{y,e}(t)|^4
 \leq3K_a^2 e^{\beta V(C+C_J^{\rm mesh})}.
 \label{eq:fixed-volume-evaluation-moment}
\end{equation}
The same bound holds for $\mu_{L,h}$. It is not a spatially uniform
bound; the later DLR passage needs its own estimate. In particular an
integrated quartic norm has not been used to bound a point evaluation
deterministically.

Define
\begin{equation}
 p_{\beta,L}(h)=\frac1V\log Z_L(h),\qquad
 P_{\beta,L}(h)=\frac{p_{\beta,L}(h)}{8\beta},\qquad
 X_L=\int_0^\beta\sum_yQ_y(\tau)\dd\tau.
 \label{eq:pressure-def}
\end{equation}
Quartic domination implies $\E_{L,0}e^{T|X_L|}<\infty$ for every finite
$T$.  Indeed, if $\mathcal Q=\int_0^\beta\sum_y|\omega_y|^4\dd\tau$,
then $|X_L|\leq(\beta V)^{3/4}\mathcal Q^{1/4}$, and the retained positive
quartic term absorbs $T|X_L|$ by Young's inequality.  Thus
$Z_L(z)=Z_L(0)\E_{L,0}e^{zX_L}$ extends to an entire function of the
complex source $z$: on $|z|\leq T$, the integrand and each of its fixed-order
derivatives have a common integrable exponential majorant.  Since
$Z_L(h)>0$ for every real $h$, it has a nonvanishing complex neighborhood
of each such point, on which a holomorphic logarithm exists.  It follows
that $p_{\beta,L}$ is real analytic on $\R$.  An entire continuation of
$\log Z_L$ would require exclusion of complex zeros and is not asserted.
For real $h$,
\begin{align}
 p_{\beta,L}'(h)&=\frac{\E_{L,h}X_L}{V}=\beta\,\E_{L,h}Q_0(0),
 &P_{\beta,L}'(h)&=\frac{\E_{L,h}Q_0(0)}8,\label{eq:source-derivatives}\\
 p_{\beta,L}''(h)&=\frac{\Var_{L,h}(X_L)}V\geq0. \nonumber
\end{align}
At $h=0$, parity $q\mapsto-q$ gives $p_{\beta,L}'(0)=0$.  A cusp can
therefore occur only after the spatial limit.

\section{Thermodynamic pressure}
\label{sec:pressure}

Let $\Lambda\subset\Z^3$ be an even rectangle.  The open Hamiltonian keeps
only bonds with both endpoints in $\Lambda$; the periodic Hamiltonian on an
even cube keeps the full positive-direction bond multiset.  The same onsite
polynomial and source are used in both cases.

Here are the uniform estimates needed for the limit; a finite upper bound
alone would not exclude a limit of minus infinity.  Put $r_-=\max(-r,0)$.
For $|h|\leq H$, scalar Young inequalities give, as quadratic forms,
\begin{equation}
 H_\Lambda^{\mathrm{op}}(h)\geq Q_\Lambda-b(H)|\Lambda|,\qquad
 Q_\Lambda=\frac g8\sum_{y,e}q_{y,e}^4,\qquad
 b(H)=\frac{8r_-^2}{g}+\frac32(4/g)^{1/3}H^{4/3}.
 \label{eq:retained-quartic}
\end{equation}
To verify the constants, retain $gx^4/8$ from each scalar quartic and use
$gx^4/16-r_-x^2/2\geq-r_-^2/g$ and
$gx^4/16-|J_e||x|\geq-(3/4)(4/g)^{1/3}|J_e|^{4/3}$, where
$J_e=h/\sqrt8$ and $\sum_e|J_e|^{4/3}=2|h|^{4/3}$.

For $\beta\in[b_0,b_1]\subset(0,\infty)$, the variance of one harmonic
reference coordinate is
\[
 k_a(\beta)=\frac{1}{2m\omega_a}\coth(\beta\omega_a/2)
 \leq K:=\frac{1}{b_0a}+\frac{b_1}{12m}.
\]
The inequality follows from
$\coth x\leq x^{-1}+x/3$ for $x>0$; for example the partial-fraction
expansion of $\coth$ bounds the remainder by
$2x\sum_{n\geq1}(\pi n)^{-2}=x/3$.
The scalar quartic and each internal edge have Gaussian means
$3k_a^2$ and $8k_a^2$, respectively.  There are eight coordinates and
twelve internal edges per site, and at most three spatial bonds per site.
The centered source has mean zero.  Consequently, for either boundary
condition,
\begin{align}
 C_J&=4|r-a|K+24cK+(6g+24\lambda)K^2,\nonumber\\
 C&=\frac{4(r-a)^2}{g}
             +\frac34(32/g)^{1/3}H^{4/3},\nonumber\\
 -8\log[2\sinh(\beta\omega_a/2)]-\beta C_J
 &\leq |\Lambda|^{-1}\log Z_\Lambda(\beta,h)
 \leq-8\log[2\sinh(\beta\omega_a/2)]+\beta C.
 \label{eq:two-sided-pressure}
\end{align}
The lower bound is Jensen's inequality in \eqref{eq:fk}; the upper bound
uses $R_{a,h}\geq-C|\Lambda|$, obtained by reserving half the quartic
term and absorbing the quadratic and source terms in the other half.
All constants are independent of $\Lambda$.  In particular the absolute
value of every normalized log trace is bounded by some $M<\infty$ on the
displayed compact parameter rectangle.

\begin{proposition}[Open-box pressure from the finite-volume representation]
For fixed $\beta>0$ and $h\in\R$, the limit
\begin{equation}
 p_\beta^{\mathrm{op}}(h)=\lim_{\Lambda\nearrow\Z^3}
 |\Lambda|^{-1}\log Z_\Lambda^{\mathrm{op}}(h)
 \label{eq:open-pressure}
\end{equation}
exists along arbitrary even rectangular exhaustions and equals the infimum
of the normalized open-box values.  It is finite, even, and convex in $h$.
\end{proposition}

\begin{proof}
Cutting a rectangle into two rectangles deletes nonnegative crossing bonds.
Min--max comparison and tensor-product trace factorization give
$F(\Lambda)=\log Z_\Lambda^{\mathrm{op}}\leq
F(\Lambda_1)+F(\Lambda_2)$.  Fix an even block $B$ with side lengths $b_i$
and tile a rectangle with even side lengths $l_i$ by copies of $B$ and
even remainder rectangles.  If $v_{\rm rem}$ is the remainder volume, then
\[
 \frac{v_{\rm rem}}{|\Lambda|}
 =1-\prod_{i=1}^3\left(1-\frac{l_i\bmod b_i}{l_i}\right)
 \leq\sum_{i=1}^3\frac{l_i\bmod b_i}{l_i}\longrightarrow0.
\]
Subadditivity and \eqref{eq:two-sided-pressure} bound $F(\Lambda)/|\Lambda|$
above by
$(1-v_{\rm rem}/|\Lambda|)F(B)/|B|+Mv_{\rm rem}/|\Lambda|$.
Therefore its limsup is at most $F(B)/|B|$ for every $B$.  Every term is
at least $\inf_B F(B)/|B|$, which is finite by the lower bound in
\eqref{eq:two-sided-pressure}.  This proves the asserted limit and infimum
formula.  Convexity and evenness pass from finite volume to this finite limit.
\end{proof}

For a periodic cube, write $B_L^{\rm seam}=H_L^{\rm per}(h)-H_L^{\rm op}(h)$
for the multiplication form given by the added seam bonds.  It is distinct
from the collective displacement Hessian $B_L$ in Section~\ref{sec:collective}.
The seam contains $3L^2$ bonds and $48L^2$ scalar
endpoint occurrences.  The endpoint Young estimate, with an arbitrary
$\eta>0$, gives the corrected bound
\begin{equation}
 0\leq B_L^{\rm seam}\leq \eta Q_L+\frac{288c^2L^2}{\eta g},\qquad
 Q_L=\frac g8\sum_{y,e}q_{y,e}^4. \label{eq:seam}
\end{equation}
For clarity, there are at most three seam endpoint occurrences at any
scalar coordinate, and
$cx^2\leq(\eta g/24)x^4+6c^2/(\eta g)$.
Thus the coefficient in \eqref{eq:seam} is $48\times6=288$.
Combining this bound with \eqref{eq:retained-quartic} on the common form
domain gives
\begin{align}
 H_L^{\rm op}&\leq H_L^{\rm per}
       \leq(1+\eta)H_L^{\rm op}+D_L^{\rm seam},\nonumber\\
 D_L^{\rm seam}&=\eta b(H)L^3+\frac{288c^2L^2}{\eta g},\nonumber\\
 e^{-\beta D_L^{\rm seam}}Z_L^{\rm op}(\beta(1+\eta),h)
 &\leq Z_L^{\rm per}(\beta,h)\leq Z_L^{\rm op}(\beta,h).
 \label{eq:seam-trace-sandwich}
\end{align}
The last line follows from min--max eigenvalue comparison and summation.

Write $f_L(\beta,h)=L^{-3}\log Z_L^{\rm op}(\beta,h)$.
On $[\beta_0/2,2\beta_1]\times[-H-1,H+1]$, choose the uniform absolute
bound $M$ supplied by \eqref{eq:two-sided-pressure}.  For fixed $h$,
$f_L$ is convex in $\beta$ (its second derivative is the Gibbs energy
variance per site).  Endpoint secants therefore bound its slopes on
$[\beta_0,3\beta_1/2]$ by $4M/\beta_0$ in absolute value.
For $0<\eta\leq1/2$ and $\beta\in[\beta_0,\beta_1]$ this gives
\[
 |f_L(\beta(1+\eta),h)-f_L(\beta,h)|
 \leq(4M\beta_1/\beta_0)\eta.
\]
Taking $\eta=L^{-1/2}$ in \eqref{eq:seam-trace-sandwich} yields
\begin{equation}
 0\leq f_L(\beta,h)-p_{\beta,L}^{\rm per}(h)
 \leq\left(\frac{4M\beta_1}{\beta_0}+\beta_1b(H)
                    +\frac{288\beta_1c^2}{g}\right)L^{-1/2}.
 \label{eq:pressure-seam-rate}
\end{equation}
Endpoint secants in $h$ also give a uniform slope bound on $[-H,H]$.
Comparing two parameter points one coordinate at a time proves common
local Lipschitz bounds.  A finite net upgrades the pointwise open-box
limit to local uniform convergence; \eqref{eq:pressure-seam-rate} transfers
it to periodic cubes.  Separate convexity suffices; joint convexity in
$(\beta,h)$ is not needed.  Thus
\begin{equation}
 p_\beta(h)=\lim_{L\to\infty,\ L\ {\rm even}}p_{\beta,L}^{\rm per}(h)
                  =p_\beta^{\rm op}(h). \label{eq:pressure-limit}
\end{equation}
The limit is locally uniform in $(\beta,h)$, finite, even, and convex in $h$.
At every differentiability point of $P_\beta(h)=p_\beta(h)/(8\beta)$,
\begin{equation}
 P_{\beta,L}'(h)\longrightarrow P_\beta'(h).
\label{eq:derivative-limit}
\end{equation}
Indeed, at fixed $\beta$ and any sufficiently small $a_1>0$, convexity gives
\[
 \frac{P_{\beta,L}(h)-P_{\beta,L}(h-a_1)}{a_1}
 \leq P_{\beta,L}'(h)
 \leq\frac{P_{\beta,L}(h+a_1)-P_{\beta,L}(h)}{a_1}.
\]
First take the volume limit in the two secants, then let $a_1\downarrow0$.
Differentiability of the limiting pressure at this fixed $h$ squeezes both
bounds to the same derivative.  No differentiability at $h=0$ is assumed.

\section{Tempered DLR states and source tangents}
\label{sec:dlr}

Let $C_\beta=C_{\rm per}([0,\beta];\R^8)$ and
$\Omega=\prod_{y\in\Z^3}C_\beta$.  For $\alpha>0$ use
\begin{equation}
 \norm{\omega}_\alpha^2=\sum_y e^{-\alpha|y|}\norm{\omega_y}_{L^2}^2,
 \qquad \Omega_t=\bigcap_{\alpha>0}\{\norm{\omega}_\alpha<\infty\},
 \label{eq:tempered-space}
\end{equation}
where $|y|$ is Euclidean distance.  Put $\alpha_k=1/k$ and enumerate the
sites as $(y_j)_{j\geq1}$.  The topology used throughout this section is
metrized by
\begin{equation}
 d_t(\omega,\eta)=\sum_{k\geq1}2^{-k}
 (1\wedge\norm{\omega-\eta}_{\alpha_k})
 +\sum_{j\geq1}2^{-j}(1\wedge\norm{\omega_{y_j}-\eta_{y_j}}_\infty).
 \label{eq:tempered-metric}
\end{equation}
The sequence $\alpha_k\downarrow0$ is cofinal for
\eqref{eq:tempered-space}.  This is a Polish space: a Cauchy sequence has
compatible limits in every weighted $L^2$ space and in every local continuous
loop space, and these limits define an element of $\Omega_t$.  Finite-support
loops with rational polygonal values at rational fractions of $\beta$ form
a countable dense set.  Weighted norms are measurable in the product
sigma-field by their finite partial sums; local sup norms are measurable by
countable time evaluations.  Thus the Borel sigma-field of this topology is
the trace of the product sigma-field.  Write $W_t$ for weak convergence of
probability measures for this topology.  Fix $0<\sigma<1/2$ and use the
periodic $C^\sigma$ norm, including the sup norm, for compactness estimates.

\subsection{Specification and the general-vector input}
\label{sec:dlr-specification}

Choose $a>0$ and the normalized oscillator Gaussian $\chi_a$.  The allocated
one-site potential and pair interaction are
\begin{align}
 V_h(q)&=\frac{r+6c-a}{2}|q|^2+\frac g4\sum_e q_e^4
 +\frac\lambda4\sum_{\{e,f\}\in E(Q_3)}(q_e-q_f)^2(q_e^2+q_f^2)-h(u,q),\\
 J_{yz}&=c\,1_{\{|y-z|=1\}},\qquad J_0=6c.
\end{align}
For a finite $\Delta$ and boundary loop $\xi$, the specification is
\begin{align}
 I_\Delta^h(\omega\mid\xi)
 &=\sum_{y\in\Delta}\int_0^\beta V_h(\omega_y)\dd\tau
 -\frac12\sum_{y,z\in\Delta}J_{yz}(\omega_y,\omega_z)_{L^2}\nonumber\\
 &\quad-\sum_{y\in\Delta,z\notin\Delta}J_{yz}(\omega_y,\xi_z)_{L^2},\\
 \pi_\Delta^h(f\mid\xi)&=Z_\Delta(h,\xi)^{-1}
 \int f(\omega_\Delta\xi_{\Delta^c})e^{-I_\Delta^h(\omega\mid\xi)}
 \prod_{y\in\Delta}\chi_a(\dd\omega_y).
 \label{eq:specification}
\end{align}
Only the negative pair interaction accompanies $V_h$; a second positive spatial
diagonal would double-count the bond energy.

Fix a source window $|h|\leq h_0$ and set
\begin{align}
 b&=\tfrac12|r+6c-a|,\qquad A=g/128,\nonumber\\
 C_0&=16b^2/g+\tfrac34 h_0^{4/3}(32/g)^{1/3},\nonumber\\
 V^+(q)&=(g/4+3\lambda)\sum_e q_e^4+b|q|^2+h_0|q|.
 \label{eq:dlr-envelope-constants}
\end{align}
Then, uniformly on this window,
\begin{equation}
 A|q|^4-C_0\leq V_h(q)\leq V^+(q).
 \label{eq:dlr-potential-envelope}
\end{equation}
Indeed, $\sum_e q_e^4\geq |q|^4/8$; absorb $b|q|^2$ against
$g|q|^4/64$ and $h_0|q|$ against $g|q|^4/128$.  The respective scalar
maxima are the two terms in $C_0$, leaving $A|q|^4$.  For the upper
bound, $(x-y)^2(x^2+y^2)\leq4(x^4+y^4)$ and the degree of $Q_3$ is
three.  All envelopes are continuous and $V_h(0)=V^+(0)=0$.

Here is the precise general-vector use of Kozitsky--Pasurek (KP)
\cite{KP}.  In their Assumption (A), (2.5), take dimension $\nu=8$,
growth exponent $r_{\rm KP}=2$, lower constants $A_V=A$, $B_V=-C_0$,
and upper envelope $V^+$.  Positive $m,a$ specify the oscillator reference.
Their interaction condition (2.6) holds with symmetric zero-diagonal $J$
and row sum $J_0=6c$; the onsite locking is not part of $J$.  Lattice
regularity (2.1) follows from
$\sum_{y\in\Z^3}(1+|y|)^{-3-\varepsilon}<\infty$ for every
$\varepsilon>0$.  For Assumption (B), use the exponential weights of
(2.42)--(2.43): $J_\alpha=6ce^\alpha<\infty$,
$J_\alpha-J_0\to0$, and the logarithmically weighted sum in (2.39) is
finite.  Their Theorem 3.1 therefore supplies a nonempty $W_t$-compact set
$\cG_t(h)$ at each fixed source.  Theorem 3.2 gives a moment bound for
this fixed-model set, not automatically for a source window or finite tori.
The topology agrees with KP (2.47)--(2.49): both require all weighted
$L^2$ norms and all local sup norms.  Replacing positive summable local
weights or $s/(1+s)$ by $1\wedge s$ leaves that topology unchanged;
$\alpha_k\downarrow0$ gives a countable cofinal family of weighted norms.
No scalar or rotation-invariant phase theorem is used.  We prove the needed
periodic and varying-source statements directly below.

\subsection{A source-window bound uniform in periodic volume}
\label{sec:dlr-periodic-moment}

For the oscillator Gaussian choose $\ell_\sigma>0$ so that
$F_\sigma:=\int e^{\ell_\sigma\norm{\omega}_{C^\sigma}^2}
\chi_a(\dd\omega)<\infty$ (KP Proposition 2.2, (2.27)).
Fix $\kappa,\theta>0$ and define
\begin{align}
 Y_-&=\int\exp\left[-\frac{J_0}{2\theta}\norm{\omega}_{L^2}^2
                   -\int_0^\beta V^+(\omega)\dd\tau\right]
                   \chi_a(\dd\omega),\nonumber\\
 D&=\kappa+J_0/(2\theta),\qquad
 C_1=\max\{0,\log F_\sigma+\beta C_0+\beta D^2/(4A)-\log Y_-\},
 \nonumber\\
 f_y(\omega)&=\ell_\sigma\norm{\omega_y}_{C^\sigma}^2
                         +\kappa\norm{\omega_y}_{L^2}^2.
 \label{eq:one-site-constants}
\end{align}
The integrand defining $Y_-$ is positive and at most one, so
$0<Y_-\leq1$.  Bilinear Young bounds give, for every neighbor boundary,
\begin{align*}
 Z_{\{y\}}(h,\xi)&\geq
 Y_-\exp\left[-\frac\theta2\sum_zJ_{yz}\norm{\xi_z}_{L^2}^2\right],\\
 \int e^{f_y-I_{\{y\}}^h}\dd\chi_a&\leq
 F_\sigma e^{\beta C_0+\beta D^2/(4A)}
 \exp\left[\frac\theta2\sum_zJ_{yz}\norm{\xi_z}_{L^2}^2\right].
\end{align*}
The second inequality uses $D s^2-A s^4\leq D^2/(4A)$ pointwise
in time.  Division, with the positive lower bound already established, yields
\begin{equation}
 \pi_{\{y\}}^h(e^{f_y}\mid\xi)
 \leq\exp\left[C_1+\theta\sum_zJ_{yz}\norm{\xi_z}_{L^2}^2\right].
 \label{eq:one-site-exponential}
\end{equation}
This is the finite-range, source-window version of the mechanism in KP
Lemma 4.1.  Its constants do not depend on the site or torus size.

Let $L\geq4$ be even.  Under the periodic law, translation symmetry makes
$M=\E_{L,h}e^{f_y}$ independent of $y$.  First $M<\infty$: in the
finite product Gaussian integral, \eqref{eq:dlr-potential-envelope} and
$|\sum_{\{y,z\}}c(\omega_y,\omega_z)_{L^2}|
\leq (J_0/2)\sum_y\norm{\omega_y}_{L^2}^2$ leave a positive quartic
at every site.  It absorbs all quadratic terms, including
$\kappa\norm{\omega_y}_{L^2}^2$; the remaining
$e^{\ell_\sigma\norm{\omega_y}_{C^\sigma}^2}$ is Gaussian integrable.
The finite normalization is positive.  This proves finiteness before any
use of a power of $M$.

Set $t=\theta J_0/\kappa<1$.  Conditional expectation,
\eqref{eq:one-site-exponential}, generalized Holder with weights
$J_{yz}/J_0$, and concavity of $x^t$ imply
\begin{align}
 M&\leq e^{C_1}\E_{L,h}
          \exp\left(\theta\sum_zJ_{yz}\norm{\omega_z}_{L^2}^2\right)
 \nonumber\\
 &\leq e^{C_1}\prod_{z:J_{yz}>0}
       (\E_{L,h}e^{tf_z})^{J_{yz}/J_0}
 =e^{C_1}\E_{L,h}e^{tf_y}\leq e^{C_1}M^t.
 \label{eq:periodic-holder-closure}
\end{align}
Choose $\theta=\kappa/(2J_0)$ and evaluate $C_1$ at this choice.
Since $M$ is finite and positive, $M\leq e^{2C_1}=:C_{\rm per}$.
In particular
\begin{equation}
 \sup_{L\geq4,\ |h|\leq h_0,y}\E_{L,h}
 \exp\left(\ell_\sigma\norm{\omega_y}_{C^\sigma}^2
                   +\kappa\norm{\omega_y}_{L^2}^2\right)
 \leq C_{\rm per}. \label{eq:periodic-moment}
\end{equation}
Extend the law by zero outside the centered cube
$[-L/2,L/2)^3\cap\Z^3$, denoting it by $\widehat\mu_{L,h}$.
Since $C_{\rm per}\geq1$, the estimate holds also at exterior sites.
In particular $\widehat\mu_{L,h}(\norm{\omega_y}_{C^\sigma}^2)
\leq C_{\rm per}/\ell_\sigma$.

\subsection{Weighted compactness with the correct tail direction}
\label{sec:dlr-weighted-compactness}

Put $\norm{\omega}_{\alpha,\sigma}^2
=\sum_y e^{-\alpha|y|}\norm{\omega_y}_{C^\sigma}^2$ and
$S(\alpha)=\sum_y e^{-\alpha|y|}<\infty$.
For a given $\varepsilon>0$, choose positive $\varepsilon_k$ with
$\sum_k\varepsilon_k\leq\varepsilon$, and radii satisfying
$R_k^2\geq C_{\rm per}S(\alpha_k)/(\ell_\sigma\varepsilon_k)$.
Markov's inequality and a union bound show that
\begin{equation}
 K_\varepsilon=\{\omega\in\Omega_t:
                  \norm{\omega}_{\alpha_k,\sigma}\leq R_k\text{ for all }k\},
 \qquad
 \inf_{L,h}\widehat\mu_{L,h}(K_\varepsilon)\geq1-\varepsilon.
 \label{eq:dlr-compact-set}
\end{equation}
This set is compact in $d_t$.  To see this, its bounds give equibounded and
equicontinuous loops at each site, hence a diagonal locally uniform subsequence
by Arzela--Ascoli.  For the weighted $L^2$ norm at $\alpha_k$, use the
\emph{stronger} bound at $\alpha_{k+1}<\alpha_k$:
\begin{equation}
 \sum_{|y|>R}e^{-\alpha_k|y|}\norm{\omega_y}_{L^2}^2
 \leq\beta e^{-(\alpha_k-\alpha_{k+1})R}R_{k+1}^2.
 \label{eq:weighted-tail-direction}
\end{equation}
This tends to zero uniformly on $K_\varepsilon$.  Finite-site convergence
therefore implies convergence in every weighted norm (apply the same tail
bound to both terms of a difference).  Lower semicontinuity of the Holder
norm under uniform convergence and Fatou's lemma put the limit back in
$K_\varepsilon$.  This proves compactness and, by Prokhorov, common
$W_t$ tightness on the entire source window.  The reverse weight direction
does not work: at $\beta=1$ a constant loop with squared amplitude $4^n$
at site $(n,0,0)$ has squared $\log4$-weighted norm one but squared
$\log2$-weighted norm $2^n$.  No spatial tightness is inferred from a
fixed-volume Gaussian bound.

\subsection{Positive finite normalizers and continuous DLR kernels}
\label{sec:dlr-kernel-continuity}

For a nonempty finite $\Delta$ of size $n$ and a compact $K\subset\Omega_t$,
write
\begin{align*}
 Q_\Delta(\omega)&=\sum_{y\in\Delta}\int_0^\beta|\omega_y|^4\dd\tau,\\
 B_K&=\tfrac12\sup_{\xi\in K}
           \sum_{y\in\Delta,z\notin\Delta}J_{yz}\norm{\xi_z}_{L^2}^2,\\
 C_{\Delta,K}&=\beta n(C_0+J_0^2/(8A))+B_K.
\end{align*}
Only finitely many boundary sites enter, so $B_K<\infty$.
The internal and exterior bilinear Young bounds allocate in total at most
$J_0/2$ to the quadratic norm of each interior loop.  Consequently
\begin{align}
 I_\Delta^h&\geq A Q_\Delta-\beta n C_0
                 -\tfrac{J_0}{2}\sum_{y\in\Delta}\norm{\omega_y}_{L^2}^2-B_K
 \nonumber\\
 &\geq\tfrac A2 Q_\Delta-C_{\Delta,K},
 \qquad |h|\leq h_0,\ \xi\in K,
 \label{eq:boundary-coercivity}
\end{align}
using $(J_0/2)s^2-(A/2)s^4\leq J_0^2/(8A)$.

For a noncircular lower bound on $Z_\Delta$, set
$R_F^2=\log(2F_\sigma)/\ell_\sigma$.  Gaussian Markov gives
$\chi_a(\norm{\omega}_\infty\leq R_F)\geq1/2$.  Define
\begin{align*}
 T_K&=\sup_{\xi\in K}\sum_{y\in\Delta,z\notin\Delta}
                       J_{yz}\norm{\xi_z}_{L^2},\\
 U_{\Delta,K}&=\beta n\{(g/4+3\lambda)R_F^4
                       +(b+J_0/2)R_F^2+h_0R_F\}
                         +\sqrt\beta R_F T_K.
\end{align*}
On the product event $\norm{\omega_y}_\infty\leq R_F$ at every interior
site, $I_\Delta^h\leq U_{\Delta,K}$, by \eqref{eq:dlr-potential-envelope}
and Cauchy--Schwarz on the boundary terms.  This event has Gaussian probability
at least $2^{-n}$.  Hence
\begin{equation}
 \inf_{|h|\leq h_0,\xi\in K}Z_\Delta(h,\xi)
 \geq z_{\Delta,K}:=2^{-n}e^{-U_{\Delta,K}}>0.
 \label{eq:dlr-normalizer-lower}
\end{equation}
Finiteness follows separately from \eqref{eq:boundary-coercivity}.

The same bounds give a useful quantitative source estimate.  Let
$X_\Delta=\sum_{y\in\Delta}\int_0^\beta(u,\omega_y)\dd\tau$.
Holder's inequality gives $|X_\Delta|\leq(\beta n)^{3/4}Q_\Delta^{1/4}$.
Because $I_\Delta^h=I_\Delta^0-hX_\Delta$, the mean value theorem and
\eqref{eq:boundary-coercivity} bound the difference of two unnormalized
weights by
\[
 |h-h'|e^{C_{\Delta,K}}(\beta n)^{3/4}
          Q_\Delta^{1/4}e^{-A Q_\Delta/2}
 \leq |h-h'|M_{\Delta,K},\qquad
 M_{\Delta,K}=e^{C_{\Delta,K}}(\beta n)^{3/4}(2Ae)^{-1/4}.
\]
The scalar maximum occurs at $Q_\Delta=1/(2A)$.  The numerator and
normalizer thus have respective Lipschitz bounds
$\norm f_\infty M_{\Delta,K}$ and $M_{\Delta,K}$.  Since the absolute
numerator is at most $\norm f_\infty$ times its own normalizer, division
gives, for every bounded Borel $f$,
\begin{equation}
 \sup_{\xi\in K}|\pi_\Delta^h(f\mid\xi)-\pi_\Delta^{h'}(f\mid\xi)|
 \leq\frac{2M_{\Delta,K}}{z_{\Delta,K}}\norm f_\infty|h-h'|.
 \label{eq:kernel-source-lipschitz}
\end{equation}

If $f\in C_b(\Omega_t)$ and $\xi_j\to\xi$ in $d_t$, replacing finitely
many interior loops is continuous, and the finite-range boundary pairings
converge in $L^2$.  The compact set consisting of this sequence and its
limit admits \eqref{eq:boundary-coercivity} and
\eqref{eq:dlr-normalizer-lower}.  Dominated convergence of the numerator
and denominator proves $\pi_\Delta^h f\in C_b(\Omega_t)$.
The kernels replace only finitely many loops, preserve $\Omega_t$, and
satisfy specification consistency by conditional integration and Fubini.
This direct finite-range Feller proof avoids assuming the desired DLR
limit in order to bound its kernel.  Since bounded continuous functions on
the Polish space determine probability measures, the equalities
$\mu(f)=\mu(\pi_\Delta^h f)$ for all such $f$ imply the full Borel DLR
equation $\mu\pi_\Delta^h=\mu$.

\subsection{Periodic limits and the two source-tangent passages}
\label{sec:dlr-source-tangents}

Fix $h$.  Once $\Delta$ and all its neighbors lie strictly within the
centered cube, the conditional law of $\widehat\mu_{L,h}$ on $\Delta$
equals \eqref{eq:specification}; no periodic seam meets an incident bond.
Thus $\widehat\mu_{L,h}(f)=\widehat\mu_{L,h}(\pi_\Delta^h f)$.
Common tightness and the Feller property pass this identity along every
converging periodic subsequence, first for $C_b(\Omega_t)$ and then for
Borel sets as above.  Each accumulation point is therefore in $\cG_t(h)$.
For every fixed local continuous cylinder test, its support and any fixed
translate eventually lie within the cube; torus symmetry then passes to
the limit.  Such tests determine the product Borel law, proving spatial
translation invariance of the limit.  The finite zero extension itself is
not spatially translation invariant.  Time-shift invariance passes too;
time shifts are continuous in $d_t$.  By lower semicontinuity and truncation,
all selected limits inherit \eqref{eq:periodic-moment}.

At a differentiability point $h$ of $P_\beta$, let $\mu_h$ be any such
accumulation point.  With $Q_0(\omega)=(u,\omega_0(0))$, define
$\operatorname{clip}_R(s)=\max(-R,\min(s,R))$.  For the periodic laws and
their selected limits,
\begin{equation}
 \E Q_0^2\leq C_{\rm per}/\ell_\sigma,\qquad
 \E|Q_0-\operatorname{clip}_R(Q_0)|
 \leq C_{\rm per}/(\ell_\sigma R).
 \label{eq:dlr-local-clipping}
\end{equation}
Abbreviate $p_L=p_{\beta,L}$ and $P_L=p_L/(8\beta)$ at this fixed
$\beta$.  For finite $L$, spatial and time translation symmetry and the actual
source $e^{hX_L}$ give
$p_L'(h)=\E X_L/V=\beta\E Q_0$ and
$P_L'(h)=\E Q_0/8$.  Section~\ref{sec:pressure} gives
$P_L'(h)\to P_\beta'(h)$ at differentiability points.  Pass bounded
clips through weak convergence and then let $R\to\infty$ in
\eqref{eq:dlr-local-clipping}; this proves
\begin{equation}
 \mu_h(Q_0)=8P_\beta'(h). \label{eq:tangent-identity}
\end{equation}
There is no extra factor $\beta$ in the normalized tangent identity.

Choose positive differentiability points $h_j\downarrow0$.  A finite convex
pressure has at most countably many nondifferentiability points.  Monotonicity
and the secant inequalities imply $P_\beta'(h_j)\to D_+P_\beta(0)$.
More explicitly, $P_\beta'(h_j)\geq D_+P_\beta(0)$; for fixed $t>0$
and $h_j<t$, convexity gives
$P_\beta'(h_j)\leq[P_\beta(t)-P_\beta(h_j)]/(t-h_j)$.
Take $j\to\infty$, then $t\downarrow0$.  The endpoint derivative is finite
because zero is interior to the finite-pressure domain.
The common source-window bound gives a subsequence
$\mu_{h_j}\Rightarrow\mu_+$ in $W_t$.
For $f\in C_b(\Omega_t)$, split the difference between the $h_j$ and
zero kernels over $K_\varepsilon$ and its complement.  Equations
\eqref{eq:dlr-compact-set} and \eqref{eq:kernel-source-lipschitz} give
\[
 |\mu_{h_j}(\pi_\Delta^{h_j}f-\pi_\Delta^0 f)|
 \leq\frac{2M_{\Delta,K_\varepsilon}}{z_{\Delta,K_\varepsilon}}
       \norm f_\infty h_j+2\norm f_\infty\varepsilon.
\]
The compact-set probability bound holds for these limits because they
inherit the moment estimate.  For fixed $\varepsilon$, take $j\to\infty$;
the zero-source Feller property passes the remaining identity.  Then let
$\varepsilon\downarrow0$.  This proves the full zero-source DLR equation.
Finally apply \eqref{eq:dlr-local-clipping} once more, now uniformly in
$j$, to pass the unbounded local expectation.  Together with
\eqref{eq:tangent-identity} this gives
\begin{equation}
 \mu_+\in\cG_t(0),\qquad
 \mu_+(Q_0)=8D_+P_\beta(0), \label{eq:source-zero-tangent}
\end{equation}
without inferring unbounded-moment convergence from weak convergence alone.
The global parity map $\Theta\omega=-\omega$ preserves $\Omega_t$ and
intertwines the zero-source specification with itself.  Hence
$\Theta_*\mu_+\in\cG_t(0)$ with opposite collective expectation.
Distinctness still requires the strictly positive cusp established later;
it does not follow merely from this construction.  Neither uniqueness nor
exhaustion of all tempered DLR states by periodic limits is asserted.  The
order is the fixed-volume loop limit, spatial accumulation at each fixed
$h$, then $h_j\downarrow0$; no arbitrary simultaneous $h=h(L)$ or
$\beta\to\infty$ passage is used.

\section{Continuous-loop FKG and collective covariance}
\label{sec:fkg}

Order loops coordinatewise in space, component, and imaginary time.  A law is
associated if $\operatorname{Cov}(F,G)\geq0$ for bounded increasing Borel
functions $F,G$.  At mesh $N$, the density is $\rho_N\propto e^{-S_N}$, where
the temporal quadratic bonds, spatial bonds, and internal $Q_3$ terms satisfy
\begin{equation}
 \partial_i\partial_j\log\rho_N
 =-\partial_i\partial_j S_N\geq0\quad(i\neq j),\qquad
 -\partial_x\partial_y\frac\lambda4(x-y)^2(x^2+y^2)
 =\frac\lambda4[(x+y)^2+5(x-y)^2]\geq0. \label{eq:log-supermodular}
\end{equation}
The first inequality concerns the log density, with the opposite sign to
the energy Hessian.  Temporal and spatial bonds contribute $m/\eps$ and
$\eps c$ per occurrence, respectively, and the internal contribution is
$\eps$ times the displayed polynomial.  Unary terms contribute zero.

Here is the finite-dimensional association argument underlying the FKG
mechanism \cite{FKG}.  First restrict a positive smooth density $\rho$
with $\partial_i\partial_j\log\rho\geq0$ to a compact product interval.
In one dimension, for independent copies $T,T'$,
\[
 \operatorname{Cov}(f(T),g(T))
 =\tfrac12\E[(f(T)-f(T'))(g(T)-g(T'))]\geq0
\]
for bounded increasing $f,g$.  Induct on dimension, conditioning on the
last coordinate $t$.  The conditional laws satisfy the induction
hypothesis.  For $t_2>t_1$ their unnormalized likelihood ratio
$R(x)=\rho(x,t_2)/\rho(x,t_1)$ is bounded and increasing, because
\[
 \partial_i\log R(x)
 =\int_{t_1}^{t_2}\partial_i\partial_t\log\rho(x,t)\dd t\geq0.
\]
Conditional association implies
$\E_{t_1}(fR)\geq\E_{t_1}f\,\E_{t_1}R$, hence
$\E_{t_2}f\geq\E_{t_1}f$.  For increasing $F(x,t)$, increase $t$
inside $F$ and then increase the conditional law; this shows that
$m_F(t)=\E_tF(x,t)$ is increasing.  The same applies to $G$, so
\[
 \operatorname{Cov}(F,G)
 =\E\operatorname{Cov}(F,G\mid t)+\operatorname{Cov}(m_F(t),m_G(t))
 \geq0.
\]
Apply the result to $\rho_N$ on growing product cubes.  Its integrability
and dominated convergence for the three bounded expectations remove the
cutoff, proving association on the full mesh space.

\subsection{Passage to bounded continuous and Borel loop tests}
\label{sec:fkg-loop-passage}

Periodic linear interpolation is order preserving.  For bounded continuous
increasing $F,G$ on $\mathsf E_L$, finite association therefore applies to
$F\circ\mathcal I_N$ and $G\circ\mathcal I_N$.  Apply the actual interacting
weak convergence \eqref{eq:weighted-loop-limit} to the three bounded
continuous functions $F,G,FG$.  Their limiting covariance is nonnegative;
$FG$ need not be increasing for this weak-convergence step.

Here is the measurable extension used below. In a Polish topological vector
space with a closed positive cone and a compatible translation-invariant
metric $d$, let $C$ be a closed upper set. If $x\leq x'$, then
$z+(x'-x)\in C$ for $z\in C$ and
$d(x',z+(x'-x))=d(x,z)$. Thus $d(x,C)$ is nonincreasing in $x$.
The bounded continuous increasing functions
$f_j(x)=\max(0,1-jd(x,C))$ decrease to $1_C$; use $f_j=0$ if $C$ is empty.
Bounded convergence transfers positive correlation to pairs of closed upper
sets.

For upper Borel sets $A,B$, choose compact subsets $K\subset A$ and
$J\subset B$ whose probabilities approximate theirs from below.  The sets
$K+E_+$ and $J+E_+$ are closed upper sets contained in $A,B$:
compactness allows extraction of the compact summand, and closedness of
the cone then passes the other summand to the limit.  Hence
\[
 \mu(A\cap B)\geq\mu((K+E_+)\cap(J+E_+))
 \geq\mu(K+E_+)\mu(J+E_+)\geq\mu(K)\mu(J).
\]
Inner regularity proves positive correlation for $A,B$. Integrating this
inequality over upper level sets proves association for nonnegative bounded
increasing Borel functions. Adding constants handles signed functions.
The finite loop space has all the stated properties with its sup metric
and pointwise positive cone. This proves continuous-loop association at
every fixed $L,h$.

For the selected spatial limits of Section~\ref{sec:dlr}, zero extension is
order preserving. Their weak convergence in the declared weighted topology
therefore passes association for bounded continuous tests. The same argument
applies along the subsequent source-tangent sequence; the Borel extension
uses the closed cone and compatible translation-invariant metric of that
Polish projective space. It does not show association for arbitrary mixtures
of DLR states. For example, the equally weighted mixture of point masses at
$(0,1)$ and $(1,0)$ has coordinate covariance $-1/4$, although each point
mass is associated.

\subsection{Coordinate products and the collective inequalities}
\label{sec:fkg-products}

At fixed volume, \eqref{eq:fixed-volume-evaluation-moment} permits clipping
of coordinate products. Along the spatial/source sequences,
\eqref{eq:periodic-moment} instead gives a common bound
$\E|\omega_{y,e}(t)|^4\leq C_4:=2C_{\rm per}/\ell_\sigma^2$, by
$z^2\leq2e^z$ with $z=\ell_\sigma\|\omega_y\|_{C^\sigma}^2$.
For two coordinates $Y,Z$ and $\operatorname{clip}_R(t)=\max(-R,\min(t,R))$,
association applies to their increasing clips. Splitting the error over
$\{|Y|>R\}$ and $\{|Z|>R\}$, Cauchy--Schwarz and the fourth-moment
Markov bound give
\begin{equation}
 \E|YZ-\operatorname{clip}_R(Y)\operatorname{clip}_R(Z)|
 \leq 2C_4/R^2.
 \label{eq:fkg-product-clipping}
\end{equation}
Clipped first moments converge as well, for example with error at most
$C_4/R^3$. This both removes the clips and permits passage of the
expectations through the stated limits.

For the finite periodic zero-source law and its parity-invariant periodic
limits, the coordinate means vanish. This parity statement is not imposed
on symmetry-breaking source-tangent states at zero source. Consequently,
on the finite zero-source law, for distinct coordinates,
\begin{equation}
 \E q_{y,e}(\tau)q_{z,f}(s)\geq0. \label{eq:coordinate-covariance}
\end{equation}
At one site, the three-regularity of $Q_3$ gives
\begin{align}
 \E D_0^{\mathrm{int}}&=3\E S_0-2\sum_{\{e,f\}\in E(Q_3)}\E(q_{0,e}q_{0,f})
 \leq3\E S_0,\\
 \E Q_0^2&=\frac18\left(\E S_0+2\sum_{e<f}\E(q_{0,e}q_{0,f})\right)
 \geq\frac18\E S_0. \label{eq:fkg-local}
\end{align}
These are expectation inequalities, not pointwise inequalities for arbitrary
vectors.

\section{Reflection positivity and the infrared bound}
\label{sec:infrared}

\subsection{A bounded positive kernel on the loop Hilbert space}
\label{sec:hilbert-kernel}

Write $B_h(q)$ for the physical onsite potential in
\eqref{eq:hamiltonian}, excluding spatial bonds.  With the normalized
oscillator Gaussian $\chi_a$ of Section~\ref{sec:dlr}, define a local finite
measure
\begin{equation}
 \nu_h(\dd\omega)=\exp\left[-\int_0^\beta
             \{B_h(\omega(\tau))-a|\omega(\tau)|^2/2\}\dd\tau\right]
             \chi_a(\dd\omega).
 \label{eq:rp-local-measure}
\end{equation}
Its density is strictly positive and bounded above by a finite constant:
the scalar quartic bounds the residual polynomial below, and the locking
term is nonnegative.  This uses the positive-difference spatial convention;
there is no additional allocated $3c|q|^2$ in this local measure.
The exact periodic law is proportional to
$\prod_y\nu_h(\dd\omega_y)\prod_{\{y,z\}\in\mathcal E_L}K(\omega_y,\omega_z)$,
where $\mathcal E_L$ consists of the nearest-neighbor bonds counted once and
\begin{equation}
 K(v,w)=\exp[-(c/2)\norm{v-w}_{\mathsf H}^2],\qquad
 \mathsf H=L^2([0,\beta];\R^8).
 \label{eq:crossing-kernel}
\end{equation}
Because $0<K\leq1$, the partition integral is finite and positive.
No infinite-dimensional Lebesgue measure is being used.

To prove positive definiteness, take finite-rank orthogonal projections
$P_n\to1$ strongly on $\mathsf H$.  On their ranges, of dimension $d_n$,
the ordinary Gaussian Fourier formula is
\[
 K_n(v,w):=e^{-c|P_n(v-w)|^2/2}
 =\int_{\R^{d_n}}e^{i k\cdot P_nv}
        \overline{e^{i k\cdot P_nw}}
        \frac{e^{-|k|^2/(2c)}}{(2\pi c)^{d_n/2}}\dd k.
\]
For every finite complex measure $\zeta$ on $\mathsf H$, Fubini therefore
gives
\begin{equation}
 \iint K_n(v,w)\zeta(\dd v)\overline{\zeta(\dd w)}\geq0.
 \label{eq:hilbert-kernel-positive}
\end{equation}
Since $|K_n|\leq1$ and $K_n\to K$ pointwise, dominated convergence for
$|\zeta|\otimes|\zeta|$ gives the same inequality for $K$.  The argument
also applies on $\mathsf H^M$ for any finite number $M$ of crossing bonds.
It does not postulate a Gaussian with identity covariance on an infinite
Hilbert space or require a positive Fourier transform of the onsite quartic.

\subsection{Direct spatial reflection positivity}
\label{sec:spatial-reflection}

For even $L\geq4$, use representatives $0,\ldots,L-1$ and reflect by
$\vartheta(y_1,y_2,y_3)=(L-1-y_1,y_2,y_3)$.
Let $\Lambda_+=\{y:y_1<L/2\}$ and $\Lambda_-=\vartheta\Lambda_+$.
The two periodic separating planes lie between sites.  Each crossing bond
joins $y\in\Lambda_+$ to $\vartheta y$, and there are $M=2L^2$ such
bonds.  Let $b(v)$ list their plus-side loops in $\mathsf H^M$.
The finite measure $\alpha_h$ on plus-half configurations is the product
of their local $\nu_h$ measures times all bonds internal to that half.
Reflection identifies the minus half with a second copy of the same measure.
The full unnormalized law is consequently
\[
 \alpha_h(\dd v)\alpha_h(\dd w)
                e^{-c\norm{b(v)-b(w)}_{\mathsf H^M}^2/2}.
\]
For a bounded complex Borel function $F$ of the plus half, push the finite
complex measure $F(v)\alpha_h(\dd v)$ through $b$ and use
\eqref{eq:hilbert-kernel-positive}.  After division by the positive
normalizer, this proves
\begin{equation}
 \E_{L,h}[F\Theta_{\vartheta}F]\geq0,\qquad
 \Theta_{\vartheta}F(\omega)=\overline{F(\vartheta\omega)}.
 \label{eq:spatial-reflection-positive}
\end{equation}
The other coordinate directions are identical.  This includes bounded
measurable tests and any spatially constant real source $h$.  It is spatial
reflection positivity, not a real-time reconstruction statement.  Gaussian
domination is invoked through the finite-dimensional theorem below, not
inferred from reflection positivity in a single unexplained step.

\subsection{The finite-dimensional FSS input and its exact source map}
\label{sec:fss-crosswalk}

The external input is Froehlich--Simon--Spencer (FSS) \cite{FSS}, Section 2,
Theorem 2.1, printed page 81, with proof on pages 82--84.  In the notation
used here, it applies to finite-dimensional vector spins on a periodic
rectangular nearest-neighbor lattice, with density proportional to
$e^{J\sum_{\{y,z\}}s_y\cdot s_z}\prod_y\varpi(\dd s_y)$, $J>0$.
The common finite prior must satisfy
$\int e^{a_1|s|^2}\varpi(\dd s)<\infty$ for every $a_1>0$; radiality is
not assumed.  For an edge field $\mathfrak h$ and its adjoint-divergence
source $B\mathfrak h$, the theorem is
\begin{equation}
 \E\exp\left(\sum_y s_y\cdot(B\mathfrak h)_y\right)
 \leq\exp\left(\frac{\norm{\mathfrak h}^2}{2J}\right).
 \label{eq:fss-theorem-input}
\end{equation}
This is a finite-dimensional input, not an imported loop theorem.  Its
constant has no dependence on spin dimension or prior moment constants.

At a fixed time mesh $\eps=\beta/N$, set
$s_y=(\sqrt\eps x_{y,k})_{k=0}^{N-1}\in\R^{8N}$.  The actual
zero-source action, with the auxiliary harmonic split canceled, is
\begin{align}
 S_{N,L,0}(x)&=\sum_y\frac m{2\eps}\sum_k|x_{y,k+1}-x_{y,k}|^2
       +\eps\sum_{y,k}B_0(x_{y,k})
       +\frac{c\eps}2\sum_{\{y,z\},k}|x_{y,k}-x_{z,k}|^2
 \nonumber\\
 &=\sum_y\mathcal V_N(s_y)-c\sum_{\{y,z\}}s_y\cdot s_z,
 \label{eq:fss-scaled-action}
\end{align}
where $\mathcal V_N$ includes the temporal term and
$\eps\sum_k\{B_0(x_k)+3c|x_k|^2\}$.  Thus $J=c$ and the same
prior $e^{-\mathcal V_N(s)}\dd s$ is used at every spatial site.  Its
normalization cancels in Gibbs expectations.  Nonnegative temporal and
locking terms and Cauchy--Schwarz on the $8N$ coordinates give
\begin{equation}
 \mathcal V_N(s)\geq\frac g{4\eps}\sum_{k,e}s_{k,e}^4
                    +\frac{r+6c}{2}|s|^2
 \geq\frac g{32\beta}|s|^4+\frac{r+6c}{2}|s|^2.
 \label{eq:fss-prior-coercivity}
\end{equation}
This proves finiteness and all quadratic exponential moments at every fixed
$N$.  The dimension and prior integrals may depend on $N$; we need only
the dimension-free source constant in \eqref{eq:fss-theorem-input}.

Here are explicit operator domains and signs.  On vertex fields define
$(Gv)(y,j)=v_{y-e_j}-v_y$, and on edge fields define
$(Ba)_y=\sum_j\{a(y+e_j,j)-a(y,j)\}$.  Summation by parts gives
$B=G^*$ and $L_{\rm sp}=G^*G=BB^*$.  This realizes the forward
differences in the FSS divergence and their backward adjoints.  The operators
act componentwise on $\R^{8N}$.  Restrict $L_{\rm sp}^{-1}$ to the
zero-sum vertex space $V_0$; its kernel on the full space is the constants.
For a real spatial field $f$ with $\sum_yf_y=0$, define
\begin{equation}
 X_L(f)=\int_0^\beta\sum_y f_yQ_y(\tau)\dd\tau,\qquad
 K_L(f)=\frac\beta{2c}\ip{f}{L_{\rm sp}^{-1}f}. \label{eq:fss-source}
\end{equation}
For real $t$, set
\begin{align}
 \eta_y(t)&=t\sqrt\eps(f_yu)_{k=0}^{N-1},\qquad
 \mathfrak h_t=G L_{\rm sp}^{-1}\eta(t),\nonumber\\
 B\mathfrak h_t&=\eta(t),\qquad
 \norm{\mathfrak h_t}^2=\ip{\eta(t)}{L_{\rm sp}^{-1}\eta(t)}
        =\beta t^2\ip f{L_{\rm sp}^{-1}f}.
 \label{eq:fss-poisson-energy}
\end{align}
The last equality uses $N\eps=\beta$ and $|u|=1$; there is no factor
eight.  The source pairing is exactly
$\sum_y s_y\cdot\eta_y(t)=tX_{N,L}(f)$, where
$X_{N,L}(f)=\eps\sum_{y,k}f_y(u,x_{y,k})$.
The vertex norm $\norm\eta^2$ cannot replace the edge Poisson norm.
Equivalently, the spatial square completion is
\begin{equation}
 \frac c2\norm{Gs}^2-\ip\eta s
 =\frac c2\norm{Gs-\mathfrak h_t/c}^2
                        -\frac1{2c}\norm{\mathfrak h_t}^2.
 \label{eq:fss-square-shift}
\end{equation}
Its bond displacement is $\mathfrak h_t/c$, not $\mathfrak h_t$.
Applying \eqref{eq:fss-theorem-input} to the exact prior yields
\begin{equation}
 \E_{N,L,0}e^{tX_{N,L}(f)}\leq e^{K_L(f)t^2}
       \quad(t\in\R).
 \label{eq:fss-mesh-mgf}
\end{equation}

\subsection{Actual loop passage and source uniform integrability}
\label{sec:fss-loop-passage}

Periodic polygonal interpolation satisfies the exact cyclic integral identity
$X_L(f)(\mathcal I_Nx)=X_{N,L}(f)$.  The functional $X_L(f)$ is
continuous for the finite-volume loop sup norm.  The interacting weak limit
\eqref{eq:weighted-loop-limit} therefore gives convergence in distribution
of these scalar variables.  Applying it to
$\min(e^{tX_L(f)},R)$ and then taking $R\uparrow\infty$ transfers
\eqref{eq:fss-mesh-mgf} without first assuming convergence of unbounded
expectations:
\begin{equation}
 \E_{L,0}\exp(tX_L(f))\leq\exp(K_L(f)t^2). \label{eq:fss-mgf}
\end{equation}

In fact the MGFs and the required derivatives converge.  For $T\geq0$,
the two signs of \eqref{eq:fss-mesh-mgf} give
\begin{equation}
 \sup_N\E_{N,L,0}e^{T|X_{N,L}(f)|}\leq2e^{K_L(f)T^2}.
 \label{eq:fss-source-ui}
\end{equation}
Using $T=2|t|$ bounds the squared $L^2$ norm of $e^{tX_{N,L}(f)}$,
which proves uniform integrability.  More generally, for every integer
$k\geq0$ and auxiliary $a_2>0$, the exponential series implies
\[
 |x|^{2k}e^{2T|x|}\leq(2k)!a_2^{-2k}e^{(2T+a_2)|x|}.
\]
Together with \eqref{eq:fss-source-ui} this bounds the $L^2$ norms of
$|X_{N,L}|^ke^{T|X_{N,L}|}$ uniformly in $N$.  Truncation and weak
convergence now pass these source moments, in particular
$\E X_{N,L}^2\to\E X_L^2$.  The parameter $a_2$ has reciprocal source-
observable units and introduces no physical assumption.  Means vanish at
zero source by parity.  All these estimates are at fixed $L,\beta,f$;
they neither control the spatial constant mode nor replace the general-source
quartic bounds or the volume-uniform DLR argument.

\subsection{Duhamel normalization and the Fourier eigenvalue}
\label{sec:infrared-duhamel}

Define the time-averaged Duhamel matrix
\begin{equation}
 C_L(y,z;\tau)=\E[Q_y(\tau)Q_z(0)],\qquad
 D_L(y,z)=\frac1\beta\int_0^\beta C_L(y,z;\tau)\dd\tau. \label{eq:duhamel-matrix}
\end{equation}
All expectations here are in the exact zero-source periodic loop law.
Its finite coordinate moments justify Fubini.  With
$\overline Q_y=\beta^{-1}\int_0^\beta Q_y(\tau)\dd\tau$, time-shift
invariance and parity give
\begin{equation}
 \E\overline Q_y\overline Q_z
 =\frac1{\beta^2}\int_0^\beta\int_0^\beta
      C_L(y,z;\tau-\sigma)\dd\tau\dd\sigma
 =D_L(y,z),\qquad
 \Var(X_L(f))=\beta^2\ip f{D_Lf}.
 \label{eq:duhamel-two-time}
\end{equation}
Thus $D_L$ is a real symmetric positive-semidefinite covariance matrix.
The exponential moment bound makes the MGF twice differentiable at zero;
its value is one and its first derivative is zero.  Subtract one from both
sides of \eqref{eq:fss-mgf}, divide by $t^2$ and take $t\to0$.  This
gives $\Var(X_L(f))\leq2K_L(f)$ and hence
\begin{equation}
 \ip{f}{D_Lf}\leq\frac1{\beta c}\ip{f}{L_{\rm sp}^{-1}f}. \label{eq:duhamel-poisson}
\end{equation}
For $p=2\pi k/L$, set $e_p(y)=V^{-1/2}e^{ip\cdot y}$ and
$E(p)=\sum_{j=1}^3(1-\cos p_j)$.  Directly from the nearest-neighbor
Laplacian,
\[
 (L_{\rm sp}e_p)(y)
 =\sum_{j=1}^3\{2-e^{ip_j}-e^{-ip_j}\}e_p(y)=2E(p)e_p(y).
\]
Spatial translation invariance diagonalizes $D_L$ in the same basis; denote
its eigenvalue by $\widehat D_L(p)$.  The real inequality
\eqref{eq:duhamel-poisson} extends to a complex field $f=a+ib$ by adding
the real inequalities for $a,b$: symmetry cancels the imaginary cross terms
in both Hermitian quadratic forms.  Thus there is no doubling of the
constant.  Applying it to $e_p$ for $p\neq0$ yields
\begin{equation}
 0\leq\widehat D_L(p)\leq\frac1{2\beta cE(p)}. \label{eq:infrared-bound}
\end{equation}
No inverse or estimate is applied to the constant mode.  This is a Fourier
or quadratic-form upper bound, not an entrywise position-space comparison.

\subsection{The singular three-dimensional sum and zero-mode subtraction}
\label{sec:infrared-singular-sum}

Define
\begin{equation}
 I_3=\frac1{(2\pi)^3}\int_{(-\pi,\pi]^3}\frac{\dd^3p}{E(p)},\qquad
 I_{3,L}=\frac1V\sum_{p\neq0}\frac1{E(p)},\qquad I_{3,L}\to I_3<\infty.
 \label{eq:i3}
\end{equation}
Here is a proof of the asserted convergence.  The elementary concavity
bound $\sin t\geq2t/\pi$ for $0\leq t\leq\pi/2$ gives
$1-\cos p_j\geq2p_j^2/\pi^2$, hence
$E(p)\geq2|p|^2/\pi^2$ on the centered Brillouin cube.  Inclusion of
$\{\norm p_\infty\leq\delta\}$ in the radius-$\sqrt3\delta$ ball
therefore gives the continuous tail bound
\begin{equation}
 \frac1{(2\pi)^3}\int_{\norm p_\infty\leq\delta}\frac{\dd^3p}{E(p)}
 \leq\frac{\sqrt3\delta}{4}\qquad(0<\delta<\pi).
 \label{eq:infrared-continuous-tail}
\end{equation}
Choose centered integer representatives $k$ for the discrete momenta.  The
shell $\norm k_\infty=n$ has at most
$(2n+1)^3-(2n-1)^3=24n^2+2$ points, and
$E(2\pi k/L)\geq8|k|^2/L^2$.  With
$M=\lfloor\delta L/(2\pi)\rfloor$, we obtain
\begin{align}
 \frac1V\sum_{0<\norm p_\infty\leq\delta}\frac1{E(p)}
 &\leq\frac1{8L}\sum_{n=1}^{M}\left(24+\frac2{n^2}\right)
 \leq\frac{13M}{4L}\leq\frac{13\delta}{8\pi}.
 \label{eq:infrared-discrete-tail}
\end{align}
An empty sum is zero.  Outside the $\delta$ cube the integrand, with its
cutoff, is bounded and Riemann integrable, so the ordinary grid sums
converge there.  First let even $L\to\infty$ for fixed $\delta$, then
let $\delta\downarrow0$ using both tail bounds.  This proves
\eqref{eq:i3} without numerical quadrature or a decimal approximation.
The controlling dimension is the spatial dimension three, not eight.

Finally the Fourier trace identity gives
\begin{equation}
 \frac{\widehat D_L(0)}V
 =D_L(0,0)-\frac1V\sum_{p\neq0}\widehat D_L(p)
 \geq D_L(0,0)-\frac{I_{3,L}}{2\beta c}.
 \label{eq:infrared-subtraction}
\end{equation}
The left side is nonnegative, but this section alone gives no strictly
positive lower bound.  A separate bound $D_L(0,0)\geq d_\beta$ would
imply $\liminf_L\widehat D_L(0)/V\geq d_\beta-I_3/(2\beta c)$.
The collective and Falk--Bruch arguments in the next section must supply
that input before any cusp or distinct-state conclusion is drawn.

\section{Collective Duhamel lower bound}
\label{sec:collective}

\subsection{Thermal energy, translated forms and the collective moment}

In this section $L,\beta$ are fixed and the source is zero.  Write
$H=H_L(0)=-(2m)^{-1}\Delta+U$, with closed form $h$, and let
$(E_i,\psi_i)$ be an orthonormal eigenbasis.  Such a basis exists because
the finite heat trace makes the heat operator compact.  Choose a constant
$C$ so that $U+C\geq0$ and $h_C=h+C$ controls the kinetic form and
$\sum_y|q_y|^4$.  The form domain in \eqref{eq:form-domain} is therefore
unchanged by this shift.  For $s\geq0$,
\begin{equation}
 s e^{-\beta s}\leq\frac{2}{e\beta}e^{-\beta s/2},\qquad
 \rho_L(H+C)<\infty,\qquad
 \rho_L\left(\sum_y|q_y|^4\right)<\infty.
 \label{eq:collective-thermal-energy}
\end{equation}
The last two statements follow from the heat trace at $\beta/2$ and
coercivity, respectively.  In particular, all multiplication polynomials
of degree at most four are absolutely integrable.  Smooth compactly
supported functions are a form core: radial cutoffs converge in $H^1$ and
the quartic weighted norm, with their gradient terms supported in escaping
annuli; mollification then converges on each compact support, where the
potential is bounded.

Let $v_{y,e}=u_e/\sqrt V$, so that $|v|=1$, and set
$H(t)=-(2m)^{-1}\Delta+U(q+tv)$ and $\Delta_t=U(q+tv)-U(q)$.
Translation preserves $H^1$ and the quartic weighted space, using
$|q_y-tv_y|^4\leq8(|q_y|^4+|tv_y|^4)$.  Thus the translated form has the
same domain.  The Hamiltonians are unitarily equivalent and $Z(t)=Z$.
For fixed $t$, $\Delta_t$ has degree at most three in $q$ and is
absolutely Gibbs integrable.  Put $p_i=e^{-\beta E_i}/Z$ and
$d_i=\ip{\psi_i}{\Delta_t\psi_i}$, interpreted as a multiplication form.
Scalar Jensen applied to the spectral probability measure of $H(t)$ in
$\psi_i$ gives
\[
 \ip{\psi_i}{e^{-\beta H(t)}\psi_i}
 \geq e^{-\beta h_t[\psi_i]}=e^{-\beta(E_i+d_i)}.
\]
Only form-domain membership is needed for the spectral first moment.
Summation and a second scalar Jensen inequality give
\begin{equation}
 1=\frac{Z(t)}Z\geq\sum_i p_i e^{-\beta d_i}
       \geq\exp\left(-\beta\sum_i p_i d_i\right).
 \label{eq:collective-jensen-chain}
\end{equation}
The first bound also makes the intermediate exponential sum finite;
$\sum_i p_i|d_i|<\infty$ follows from the polynomial bound.  Consequently
\begin{equation}
 \rho_L\big(U(q+tv)-U(q)\big)\geq0. \label{eq:jensen-shift}
\end{equation}
The finite polynomial $t\mapsto\rho_L(\Delta_t)$ has a minimum at zero.
Its second derivative is nonnegative.  This differentiates finitely many
integrable coefficients, not the heat trace under an unbounded cubic
perturbation, and uses no operator monotonicity of the exponential.
Spatial differences and internal differences are unchanged by the common
displacement.  The mass term contributes $r$, the scalar quartic contributes
$3g\sum_y S_y/(8V)$, and each locking edge contributes
$\lambda(q_{y,e}-q_{y,f})^2/(8V)$: the second derivative of
$q_{y,e}^2+q_{y,f}^2$ is $4/(8V)$.  Hence
\begin{equation}
 B_L:=\left.\partial_t^2U(q+tv)\right|_{t=0}
 =r+\frac{3g}{8V}\sum_yS_y
       +\frac\lambda{8V}\sum_yD_y^{\mathrm{int}},\qquad
 \rho_L(B_L)\geq0. \label{eq:collective-hessian}
\end{equation}
For $\Pi_0=-i\hbar v\cdot\nabla=V^{-1/2}\sum_y(u,p_y)$, direct
differentiation on the smooth compact core gives
\begin{equation}
 [\Pi_0,[H_L(0),\Pi_0]]=\hbar^2B_L. \label{eq:double-commutator}
\end{equation}
The right side is a quadratic, form-bounded multiplier and defines the
continuous form extension of this core identity.  No Gibbs trace of a
product of three unbounded operators is used: positivity of its expectation
was established separately in \eqref{eq:collective-jensen-chain}.
The $\hbar^2$ belongs to the momentum commutator, not the displacement
Hessian.  In particular, \eqref{eq:collective-hessian} is not pointwise
positivity of $B_L$.

Spatial translation invariance gives the first inequality below.  At zero
source parity makes each coordinate mean zero, and the actual loop FKG
and clipping argument in \eqref{eq:fkg-local} makes all off-diagonal
coordinate products nonnegative in expectation.  Therefore
\begin{equation}
 -r\leq\frac{3g}{8}\rho_L(S_0)+\frac\lambda8\rho_L(D_0^{\mathrm{int}}),
 \qquad \rho_L(Q_0^2)\geq\frac18\rho_L(S_0). \label{eq:moment-chain}
\end{equation}
Indeed, the second inequality follows by expanding the square of
$Q_0=\sum_e q_{0,e}/\sqrt8$.  The three-regular internal graph gives
\begin{equation}
 \rho_L(D_0^{\mathrm{int}})
 =3\rho_L(S_0)-2\sum_{\{e,f\}\in E(Q_3)}\rho_L(q_{0,e}q_{0,f})
 \leq3\rho_L(S_0).
 \label{eq:collective-graph-expectation}
\end{equation}
These are Gibbs expectation inequalities, not pointwise graph bounds.
Together, for $r<0$, they imply
\begin{equation}
 \rho_L(Q_0^2)\geq\theta_Q,\qquad
 \theta_Q=-\frac r{3(g+\lambda)}>0. \label{eq:theta-q}
\end{equation}

\subsection{A finite spectral proof of the Falk--Bruch inequality}
\label{sec:finite-falk-bruch}

We give the scalar operator inequality used here explicitly; it is the
standard Falk--Bruch inequality, not a new model-specific result.  Its
normalization agrees with \cite[Proposition 3.18, (3.65), (3.68)]{KKK}.
For a bounded self-adjoint $A$, define
\begin{align}
 g(A)&=\sum_{i,j}p_i|A_{ij}|^2
       =\frac12\sum_{i,j}(p_i+p_j)|A_{ij}|^2,\nonumber\\
 b(A)&=\sum_{i,j}\mathcal L(p_i,p_j)|A_{ij}|^2,\qquad
 \mathcal L(a,b)=\int_0^1 a^{1-t}b^t\dd t
 =\begin{cases}(a-b)/(\log a-\log b),&a\ne b,\\a,&a=b.
 \end{cases} \label{eq:collective-log-mean}
\end{align}
The integral and weighted arithmetic--geometric mean imply
$0<\mathcal L(a,b)\leq(a+b)/2$.  Thus $0\leq b(A)\leq g(A)$, and
polarization gives a positive form with Cauchy--Schwarz.  Spectral expansion
and Tonelli identify
\begin{equation}
 b(A)=\frac1{\beta Z}\int_0^\beta
 \Tr\left(e^{-(\beta-\tau)H}A e^{-\tau H}A\right)\dd\tau.
 \label{eq:collective-duhamel-spectral}
\end{equation}
This trace is the nonnegative spectral sum, equivalently the squared
Hilbert--Schmidt norm of
$e^{-\tau H/2}A e^{-(\beta-\tau)H/2}$; it is finite, including at
the endpoints.  For $0<\tau<\beta$, the elementary bound
$p_i^{1-\tau/\beta}p_j^{\tau/\beta}\leq
(1-\tau/\beta)p_i+(\tau/\beta)p_j$ suffices.  No separately defined
positive-time exponential $e^{+\tau H}$ is needed.  For bounded
multipliers this is also the finite-volume loop two-time integral.

In a finite-dimensional Gibbs system set
\begin{equation}
 c(A)=\beta\sum_{i,j}(E_j-E_i)(p_i-p_j)|A_{ij}|^2
      =\rho([A,[\beta H,A]])\geq0.
 \label{eq:finite-spectral-c}
\end{equation}
Let $f(0)=1$ and, for $k>0$, let $x>0$ be the unique solution of
$x\tanh x=k$ and define $f(k)=\tanh x/x$.  The function $x\tanh x$
increases continuously from zero to infinity, so $f$ is well-defined and
continuous on $[0,\infty)$.

For completeness, the scalar function
$\Phi(u)=\sqrt u\coth\sqrt u$, with $\Phi(0)=1$, is increasing and
concave.  For $x>0$ differentiation gives
\begin{align}
 \Phi'(x^2)&=\frac{\sinh x\cosh x-x}{2x\sinh^2x}>0,\nonumber\\
 \Phi''(x^2)&=\frac{\cosh x}{4x^3\sinh^3x}
       \left(2x^2-x\tanh x-\sinh^2x\right)<0.
 \label{eq:falk-phi-concavity}
\end{align}
The first numerator starts at zero and has derivative $2\sinh^2x>0$.
For the second sign, $0\leq\tanh x\leq x$ and
$(\tanh x-x+x^3/3)'=x^2-\tanh^2x\geq0$ imply
$\tanh x\geq x-x^3/3$.  Taylor's formula with nonnegative remainder
gives $\sinh x\geq x+x^3/6$.  Consequently
$\sinh^2x+x\tanh x\geq2x^2+x^6/36$.  Continuity at zero extends
the concavity to the closed half-line.

For a nonzero self-adjoint finite matrix $A$, $b=b(A)>0$.  Put
$x_{ij}=|\log p_i-\log p_j|/2$ and
$w_{ij}=\mathcal L(p_i,p_j)|A_{ij}|^2/b$.  These weights sum to one,
including any degenerate-energy pairs with $x_{ij}=0$.  The exact identities
\begin{equation}
 \frac{g}{b}=\sum_{i,j}w_{ij}\Phi(x_{ij}^2),\qquad
 \frac{c}{4b}=\sum_{i,j}w_{ij}x_{ij}^2
 \label{eq:falk-spectral-weights}
\end{equation}
follow from the arithmetic/logarithmic mean ratio $x\coth x$ and
$\beta(E_j-E_i)=\log p_i-\log p_j$.  Jensen now gives
$g/b\leq\Phi(c/(4b))$.  If $c=0$, every active pair has zero energy
gap and $g=b$.  If $c>0$, set $z=\sqrt{c/(4b)}$ to obtain
$g\leq c/(4z\tanh z)$.  If $x\tanh x=c/(4g)$, monotonicity gives
$x\geq z$, and hence
\begin{equation}
 b(A)\geq g(A)f\left(\frac{c(A)}{4g(A)}\right).
 \label{eq:finite-falk-bruch}
\end{equation}
For $A=0$ interpret the right side as zero.  This proves the precise finite
inequality needed below without importing a radial moment or phase theorem.

\subsection{Bounded coordinates, energy-weighted sums and cutoff removal}

Return to the fixed finite-volume Schr\"odinger Hamiltonian.  Let
$A_R=R\tanh(Q_0/R)$ for $R>0$.  At fixed $R$, both $A_R$ and $A_R^2$
are bounded multipliers with bounded first derivatives and hence preserve
the closed form domain.  Since $\norm{u}=1$,
$|\nabla A_R|^2=\operatorname{sech}^4(Q_0/R)\leq1$.
On the smooth form core, the potential and mixed gradient terms cancel in
\[
 h[A_R\psi]-\operatorname{Re}h(\psi,A_R^2\psi)
 =\frac1{2m}\int |\nabla A_R|^2|\psi|^2\dd q.
\]
Form density passes the identity to the domain.  For an eigenvector,
$A_R^2\psi_i$ is an admissible form test, so the eigenvector form equation
gives
\begin{equation}
 h[A_R\psi_i]-E_i\norm{A_R\psi_i}^2
 =\frac1{2m}\int \operatorname{sech}^4(Q_0/R)|\psi_i|^2\dd q.
 \label{eq:form-identity}
\end{equation}
Before regrouping any energy-weighted sums, observe that
$h_C[A_R\psi]\leq C_R(h_C[\psi]+\norm{\psi}^2)$, by the product
rule, $U+C\geq0$, and the multiplier bounds.  Thus
\begin{align}
 \sum_{i,j}p_i(E_j+C)|(A_R)_{ij}|^2
  &=\sum_i p_i h_C[A_R\psi_i]<\infty,\nonumber\\
 \sum_{i,j}p_i(E_i+C)|(A_R)_{ij}|^2
  &\leq R^2\rho_L(H+C)<\infty.
 \label{eq:collective-absolute-energy-sums}
\end{align}
Here self-adjointness allows the row/column interchange for the first
identity.  Since $E_i+C\geq0$, these bounds control
$\sum_{i,j}p_i|E_j-E_i||(A_R)_{ij}|^2$; the corresponding $p_j$ sum
is equal after interchanging indices.  Regrouping is therefore legitimate,
and \eqref{eq:form-identity} gives the finite scalar
\begin{align}
 c_R&:=\beta\sum_{i,j}(E_j-E_i)(p_i-p_j)|(A_R)_{ij}|^2\nonumber\\
 &=2\beta\sum_i p_i\big(h[A_R\psi_i]-E_i\norm{A_R\psi_i}^2\big)
 =\frac\beta m\rho_L\big(\operatorname{sech}^4(Q_0/R)\big),
 \quad 0\leq c_R\leq\frac\beta m.
 \label{eq:collective-local-c}
\end{align}
Every paired spectral summand is nonnegative.  This is the form
interpretation of $\rho_L([A_R,[\beta H,A_R]])$, not a formal cyclic trace
calculation.  The scalar $c_R$ is the expectation of the displayed bounded
multiplier; the two objects are not identified.

To apply \eqref{eq:finite-falk-bruch}, let $P_M$ project onto the first
$M$ eigenvectors, $H_M=\operatorname{diag}(E_1,\ldots,E_M)$,
$A_{R,M}=P_M A_R P_M$, $Z_M=\sum_{i\leq M}e^{-\beta E_i}$ and
$q_M=Z_M/Z$.  The finite Gibbs probabilities are $p_i/q_M$.
Let $g_{R,M},b_{R,M},c_{R,M}$ be the three finite expressions.  Homogeneity
of the logarithmic mean gives
\begin{align}
 q_M g_{R,M}&=\sum_{i,j\leq M}p_i|(A_R)_{ij}|^2,\nonumber\\
 q_M b_{R,M}&=\sum_{i,j\leq M}\mathcal L(p_i,p_j)|(A_R)_{ij}|^2,\nonumber\\
 q_M c_{R,M}&=\beta\sum_{i,j\leq M}(E_j-E_i)(p_i-p_j)|(A_R)_{ij}|^2.
 \label{eq:collective-spectral-cutoff}
\end{align}
The unnormalized right sides increase to $g(A_R),b(A_R),c_R$ and
$q_M\to1$.  Thus all three normalized quantities converge at fixed
$R,L,\beta$.  Since $A_R$ is nonzero and the Gibbs state is faithful,
$g(A_R)>0$ and the finite $g_{R,M}$ are eventually positive.  Continuity
of $f$ passes \eqref{eq:finite-falk-bruch} to
\begin{equation}
 b(A_R)\geq g(A_R)f\left(\frac{c_R}{4g(A_R)}\right).
 \label{eq:collective-bounded-falk}
\end{equation}
We do not assert monotonicity of the normalized finite expressions,
operator-norm convergence of commutators, or an exact canonical commutation
relation in a finite matrix algebra.

Only now let $R\to\infty$, still at fixed $L,\beta$.  Since
$|A_R|\leq|Q_0|$, $A_R\to Q_0$ and $\rho_L(Q_0^2)<\infty$,
dominated convergence gives
\begin{equation}
 g(A_R)\longrightarrow s_L:=\rho_L(Q_0^2),\qquad
 g(A_R-Q_0)\longrightarrow0,\qquad c_R\longrightarrow\frac\beta m.
 \label{eq:collective-coordinate-cutoff}
\end{equation}
For the last limit use bounded convergence in
\eqref{eq:collective-local-c}.  The bound $b\leq g$ and Cauchy--Schwarz
extend $b$ to the thermal $L^2$ completion and give
$|\sqrt{b(A_R)}-\sqrt{b(Q_0)}|\leq\sqrt{g(A_R-Q_0)}$.
There is also direct loop identification: time invariance and
Cauchy--Schwarz bound the difference of the two-time products uniformly in
$\tau$ by
\begin{equation}
 \E_{L,0}|A_R(\tau)A_R(0)-Q_0(\tau)Q_0(0)|
 \leq\big(\sqrt{g(A_R)}+\sqrt{s_L}\big)\sqrt{g(A_R-Q_0)}.
 \label{eq:collective-loop-clipping}
\end{equation}
Here $A_R(\tau)$ denotes the clipped loop coordinate.  Integration proves
$b(Q_0)=D_L(0,0)$ without positing an unbounded operator trace.  Parity
makes both odd means zero, so raw and connected forms agree.  Therefore
\begin{equation}
 D_L(0,0)\geq s_L f\left(\frac\beta{4ms_L}\right),\qquad
 s_L=\rho_L(Q_0^2)\geq\theta_Q>0. \label{eq:duhamel-lower-general}
\end{equation}
The local coefficient is $\beta/m=\beta\hbar^2/\chi$: it contains the
$\beta$ from $\beta H$ and is distinct from the global momentum identity
\eqref{eq:double-commutator}.  To substitute the moment lower bound, fix
$k>0$ and write $x_s\tanh x_s=k/s$.  Then
$s f(k/s)=k/x_s^2$.  Increasing $s$ decreases $x_s$, so this expression
increases.  Using $k=\beta/(4m)$ yields
\begin{equation}
 D_L(0,0)\geq d_\beta:=\theta_Q\frac{\tanh x_\beta}{x_\beta},\qquad
 x_\beta\tanh x_\beta=\frac\beta{4m\theta_Q}. \label{eq:d-beta}
\end{equation}
All domain arguments were at fixed volume: finite spectral cutoff first,
$M\to\infty$ at fixed $R$, then $R\to\infty$.  The final scalar bound
has no $V$ and can now be combined with the infrared sum before taking the
spatial limit.  It constructs no common infinite-volume unbounded operator,
nor a volume-uniform operator core.

\section{Strict cusp and two DLR states}
\label{sec:cusp}

\subsection{The two proof branches and their common normalization}
\label{sec:composition-dictionary}

Fix $m,c,g,\lambda>0$ and $r<0$, and keep one $\beta>0$ throughout
the spatial and source limits.  The zero-source fluctuation estimate and
the positive-source DLR construction use the same physical Hamiltonian,
but different sequences of laws.  Their common dictionary is
\begin{equation}
 \frac{Z_L(h)}{Z_L(0)}=\E_{L,0}e^{hX_L},\qquad
 Y_L:=\frac{X_L}{V},\qquad
 \E_{L,0}Y_L^2=\beta^2\Pi_L,\qquad
 P_\beta=\frac{p_\beta}{8\beta}.
 \label{eq:composition-source-dictionary}
\end{equation}
Here $X_L$ already contains the time integral in \eqref{eq:pressure-def}.
In particular, the source in the exponential is $h$, not $\beta h$.
The moment estimate below is in $\mu_{L,0}$; the tangent states later are
limits of $\mu_{L,h_j}$ with $h_j$ held fixed during each volume limit.
No zero-source periodic accumulation is asserted to be a positive phase.

Time invariance, \eqref{eq:duhamel-two-time}, and the Fourier trace identity
give
\begin{align}
 \Pi_L&:=\E_{L,0}\left[\left(\frac{X_L}{\beta V}\right)^2\right]
 =\frac1{V^2}\sum_{y,z}D_L(y,z)
 =\frac{\widehat D_L(0)}V\nonumber\\
 &=D_L(0,0)-\frac1V\sum_{p\neq0}\widehat D_L(p). \label{eq:zero-mode}
\end{align}
Thus the positive local bound is used only after subtracting every
nonzero mode.  Combining \eqref{eq:infrared-bound}, \eqref{eq:i3}, and
\eqref{eq:d-beta} first gives
$\Pi_L\geq d_\beta-I_{3,L}/(2\beta c)$ and then
\begin{equation}
 \liminf_{L\to\infty}\Pi_L\geq
 \delta_\beta:=d_\beta-\frac{I_3}{2\beta c}. \label{eq:delta}
\end{equation}
Define
\begin{equation}
 A_0=8cm\theta_Q^2. \label{eq:a0}
\end{equation}
Because $\beta=4m\theta_Qx_\beta\tanh x_\beta$,
\begin{equation}
 2\beta c\,\delta_\beta=A_0\tanh^2x_\beta-I_3. \label{eq:threshold-algebra}
\end{equation}
If $A_0>I_3$, set
\begin{equation}
 \rho_*=\sqrt{I_3/A_0},\qquad x_*=\operatorname{artanh}(\rho_*),\qquad
 \beta_*=4m\theta_Qx_*\rho_*. \label{eq:beta-star}
\end{equation}
Then $\beta>\beta_*$ implies $\delta_\beta>0$.  At equality, or when
$A_0\leq I_3$, the lower bound is nonpositive; this is not a phase-absence
statement.
Indeed $x\tanh x$ and $\tanh^2x$ are strictly increasing for $x>0$,
so for $A_0>I_3$ the condition $\delta_\beta>0$ is equivalent to
$\beta>\beta_*$.  Also $0<I_3<\infty$: positivity follows from the
positive integrand away from the origin, and finiteness from
\eqref{eq:infrared-continuous-tail}.  The sufficient set is nonempty.
For fixed $c,g,\lambda>0$ and $r<0$, choose
$m>I_3/(8c\theta_Q^2)$ and then $\beta>\beta_*$.
This is a choice of finite parameters, not a zero-temperature or infinite-mass
limit.  The threshold is sufficient and is not identified with the exact
critical temperature.

\subsection{A squared-tail estimate at a nondifferentiable pressure}

We use the following elementary endpoint lemma rather than differentiating
finite-volume second moments through the cusp.

\begin{lemma}[Even-pressure squared-tail bound]
\label{lem:griffiths-endpoint}
Let $V_n>0$ tend to infinity and let $X_n$ have finite moment generating
functions on a common interval $(-h_0,h_0)$, $h_0>0$.  Suppose
\[
 F_n(h)=V_n^{-1}\log\E e^{hX_n}\longrightarrow F(h)
\]
pointwise on that interval, where $F$ is finite, even, and $F(0)=0$.
Then $F$ is convex, its right derivative $\ell=F'_+(0)$ is finite and
nonnegative, and
\begin{equation}
 \limsup_n\E[(X_n/V_n)^2]\leq\ell^2. \label{eq:griffiths-tail}
\end{equation}
\end{lemma}

\begin{proof}
Convexity passes from the log MGFs to their finite limit. Evenness makes
zero a minimum, and finiteness on a two-sided interval bounds the one-sided
slopes there.  Fix $\eps>0$ and choose a fixed $h\in(0,h_0)$ with
$F(h)/h\leq\ell+\eps/4$.
Pointwise convergence and evenness give, for large $n$,
$\max(F_n(h),F_n(-h))\leq h(\ell+\eps/2)$.  Chernoff's bound for
$Y_n=X_n/V_n$ gives
\[
 \Pr(|Y_n|>z)\leq2e^{-V_nh(z-\ell-\eps/2)}
 \quad(z\geq\ell+\eps).
\]
For $R=\ell+\eps$, the exact layer-cake identity and Tonelli give
\begin{align}
 \E[Y_n^2 1_{\{|Y_n|>R\}}]
 &=R^2\Pr(|Y_n|>R)+\int_R^\infty2z\Pr(|Y_n|>z)\dd z\nonumber\\
 &\leq2e^{-V_nh\eps/2}\left[R^2+\frac{2R}{V_nh}
                         +\frac2{(V_nh)^2}\right]\longrightarrow0.
 \label{eq:griffiths-squared-tail}
\end{align}
The central contribution is at most $R^2$.  Take the limsup and then
$\eps\downarrow0$.
\end{proof}

This is the needed second-moment specialization of
\cite[Proposition 3.9, (3.23)--(3.24)]{KKK}, proved here using only a
neighborhood of zero.  It uses fixed nonzero test sources before taking
$V_n\to\infty$, and controls a squared tail rather than merely weak
convergence or tail probabilities.  No derivative at zero is interchanged
with the volume limit.

Apply the lemma to the scalar pushforward of $\mu_{L,0}$ by $X_L$.
Equation~\eqref{eq:composition-source-dictionary} gives
$F_L(h)=p_{\beta,L}(h)-p_{\beta,L}(0)$, which converges to the finite,
even $F(h)=p_\beta(h)-p_\beta(0)$ by \eqref{eq:pressure-limit}.
For $\delta_\beta>0$, the two independently justified bounds give
\begin{equation}
 \beta^2\delta_\beta
 \leq\liminf_L\E Y_L^2
 \leq\limsup_L\E Y_L^2
 \leq[p_\beta'(0+)]^2.
 \label{eq:composition-liminf-limsup}
\end{equation}
The right slope is nonnegative and finite, so taking its positive square
root and dividing by $8\beta$ yields
\begin{equation}
 p_\beta'(0+)\geq\beta\sqrt{\delta_\beta},\qquad
 D_+P_\beta(0)=\frac{p_\beta'(0+)}{8\beta}
 \geq\frac{\sqrt{\delta_\beta}}8>0. \label{eq:cusp}
\end{equation}
Evenness gives $D_-P_\beta(0)=-D_+P_\beta(0)$, so the pressure has a strict
collective-source cusp in the sufficient regime.
At every finite volume, in contrast,
$p_{\beta,L}'(0)=0$ by parity.  The derivative convergence
\eqref{eq:derivative-limit} applies only at differentiability points of
the limiting pressure and makes no assertion at this cusp.

\subsection{Source-selected DLR states and a bounded witness of distinctness}

Choose positive differentiability points $h_j\downarrow0$.  Equations
\eqref{eq:source-zero-tangent} and \eqref{eq:cusp} give a zero-source DLR
state $\mu_+$ with
\begin{equation}
 \mu_+(Q_0)=8D_+P_\beta(0)\geq\sqrt{\delta_\beta}>0. \label{eq:mu-plus}
\end{equation}
For clarity, this construction uses both passages in
Section~\ref{sec:dlr-source-tangents}: at each fixed $h_j$, take a spatial
accumulation and identify its clipped local expectation with
$8P_\beta'(h_j)$; then take a source-to-zero subsequence using common
$W_t$ tightness, compact-uniform kernel continuity and the second clipping
passage.  A cusp by itself would not establish the DLR property.

Global parity $\Theta\omega=-\omega$ is a continuous involution of
$\Omega_t$.  The Gaussian reference is symmetric and
$I_\Delta^0(-\omega\mid-\xi)=I_\Delta^0(\omega\mid\xi)$, so change of
variables in the finite specification gives, for bounded Borel $f$,
\begin{equation}
 \pi_\Delta^0(f\mid\Theta\xi)
 =\pi_\Delta^0(f\circ\Theta\mid\xi).
 \label{eq:parity-specification-intertwining}
\end{equation}
Consequently
$(\Theta_*\mu_+)\pi_\Delta^0(f)
=\mu_+\pi_\Delta^0(f\circ\Theta)=\mu_+(f\circ\Theta)$.
Thus $\mu_-:=\Theta_*\mu_+$ is also a tempered DLR
state and
\begin{equation}
 \mu_-(Q_0)=-\mu_+(Q_0)<0. \label{eq:mu-minus}
\end{equation}
They are distinct because their finite local expectations differ.  One can
also exhibit a bounded continuous cylinder witness.  Let
$a_\beta=\sqrt{\delta_\beta}$ and $C_2=C_{\rm per}/\ell_\sigma$ from
\eqref{eq:dlr-local-clipping}.  Parity preserves that second-moment bound.
For $f_R=\operatorname{clip}_R(Q_0)$, the two clipping errors give
\begin{equation}
 \mu_+(f_R)-\mu_-(f_R)
 \geq2a_\beta-\frac{2C_2}{R}>a_\beta>0
 \quad\text{whenever } R>\frac{2C_2}{a_\beta}.
 \label{eq:bounded-phase-witness}
\end{equation}
This uses a finite $R$ after the states have been constructed.  They are
not asserted to be extremal, pure, clustering, or the only phases.

\begin{theorem}[Conditional Q3LOCK phase-coexistence package]
\label{thm:q3lock-composition}
Fix the Hamiltonian \eqref{eq:hamiltonian}, with $m,c,g,\lambda>0$ and
$r\in\R$, on
even periodic cubes in the fixed lattice $\Z^3$ and use the pressure and
source normalization \eqref{eq:pressure-def}.
Assume the finite-volume mesh/loop identification, the KP general-vector
Feller and compactness hypotheses, the selected-law FKG passage, the finite
FSS inequality, the collective form/Falk--Bruch passage, and the
source-window DLR continuity described in Sections~\ref{sec:finite}--\ref{sec:collective},
with the displayed uniformities and limit orders.  These are upstream
analytic inputs, not an assumption of a cusp or state multiplicity.
The pressure limit is finite, even and convex for every real source, and
the tempered DLR set at each fixed source is nonempty and $W_t$-compact.
If also $r<0$, $A_0>I_3$, and $\beta>\beta_*$, then the pressure has the
strict cusp \eqref{eq:cusp}; moreover there are at least two distinct
parity-related tempered Euclidean DLR states at $h=0$ satisfying
\eqref{eq:mu-plus}--\eqref{eq:mu-minus}.
\end{theorem}

\begin{proof}
The finite loop/trace and source identities are established in
Section~\ref{sec:finite}.  The two-sided pressure bounds, open-box limit
and periodic seam estimate in Section~\ref{sec:pressure} yield the common
finite pressure; the general-vector input in
Section~\ref{sec:dlr-specification} gives fixed-source DLR existence and
compactness.  At zero source, actual-loop FKG in Section~\ref{sec:fkg}
enters the collective moment bound, while Section~\ref{sec:collective}
turns it into $D_L(0,0)\geq d_\beta$.  The nonzero-mode and singular-sum
arguments in Section~\ref{sec:infrared} therefore yield
\eqref{eq:delta}.  The strict threshold algebra and
Lemma~\ref{lem:griffiths-endpoint}, with the exact scalar source
pushforward, give \eqref{eq:composition-liminf-limsup} and the cusp.
The separate source-window DLR/tangent construction supplies $\mu_+$;
\eqref{eq:parity-specification-intertwining} supplies $\mu_-$, and
\eqref{eq:bounded-phase-witness} distinguishes them.  Each implication uses
the same fixed $\beta$ and physical coefficients.  No simultaneous
$h=h(L)$ limit or information about all DLR mixtures is required.
\end{proof}

\begin{remark}[Status of the theorem]
The theorem is the target conditional composition of the registered proof
blocks. Detailed arguments and a source-to-conclusion composition are now
present, but their presence and finite diagnostics do not constitute signed
independent acceptance. The current repository tier is T0 and the theorem must be
reopened if any source convention, limit order, domain condition, or
normalization fails independent audit.
\end{remark}

\section{Proof dependency and literature crosswalk}
\label{sec:crosswalk}

Table~\ref{tab:crosswalk} separates reused standard machinery from the
model-specific interfaces.  Exact bibliographic details and source hashes are
in the companion \texttt{literature-crosswalk.md}.

\begin{table}[ht]
\centering
\small
\begin{tabular}{p{0.23\textwidth}p{0.31\textwidth}p{0.34\textwidth}}
\toprule
Source or input & Role & Boundary in this paper\\
\midrule
Simon, arXiv:math-ph/9907022v1 & Free-bridge formula, Theorem 1.1 & Harmonic reweighting and absolute normalization are derived here\\
Kozitsky--Pasurek, arXiv:math-ph/0609045v1 & General-vector DLR existence/compactness and Gaussian Fernique input & Feller/source-window passages are written here; no scalar phase theorem is imported\\
Froehlich--Simon--Spencer (1976) & Finite-dimensional Poisson-shift infrared inequality & Applied at fixed time mesh after an explicit $8N$-component source map; loop passage is separate\\
Fortuin--Kasteleyn--Ginibre (1971) & Finite partially ordered association mechanism & Continuous-loop association is proved by mesh induction and weak-limit passage\\
Kargol--Kondratiev--Kozitsky, arXiv:0710.2303v1 & Endpoint Griffiths and Falk--Bruch comparators & The endpoint specialization and finite spectral inequality are proved here; no radial theorem is imported\\
EXP-001587/588/589 & Corrected residual, finite pressure, and open/periodic pressure limit & Internal model-specific inputs, still awaiting signed audit\\
EXP-001591/593/595/1598 & DLR tangent, FKG, infrared, and collective lower-bound blocks & Conditional composition; finite diagnostics are not proof certification\\
\bottomrule
\end{tabular}
\caption{Proof dependency and source boundary.}
\label{tab:crosswalk}
\end{table}

FSS Theorem 2.1 already permits nonradial priors: lack of internal rotation
symmetry is not a new infrared theorem. Likewise the functional form of our
sufficient condition is standard. With $f(x\tanh x)=\tanh x/x$,
\begin{equation}
 d_\beta=\theta_Q f(\beta/(4m\theta_Q)),\qquad
 2\beta c\theta_Q f(\beta/(4m\theta_Q))>I_3.
 \label{eq:standard-threshold-comparison}
\end{equation}
This is the same scalar threshold mechanism as KKK \cite{KKK},
Theorem 3.20, (3.71)--(3.75), after replacing its model-dependent
lower-bound parameter by $\theta_Q$. This substitution compares formulas;
it does not validate that parameter for Q3LOCK. That requires
Section~\ref{sec:collective}.
The coefficient eight in $A_0$ arises here from the factor two in the
infrared subtraction and the factor four in the Falk--Bruch argument,
not from counting the eight internal coordinates.

The precise obstruction to directly importing KKK Theorems 3.20--3.21 is
their radial-potential assumption, not the spatial lattice. Let $W$ denote
the quartic part of $B_0$, and let $e_0$ be a coordinate unit vector.
The vectors $e_0$ and $u$ both have norm one, but
\begin{equation}
 W(e_0)=\frac{g+3\lambda}{4},\qquad W(u)=\frac g{32},\qquad
 W(e_0)-W(u)=\frac{7g}{32}+\frac{3\lambda}{4}>0.
 \label{eq:nonradial-witness}
\end{equation}
Only the three edges incident to $e_0$ contribute in the first evaluation;
all locking differences vanish in the second. Thus even an orthogonal
change of coordinates cannot make this potential radial.

Flattening $(y,e)$ into scalar sites does not give the usual scalar model
with only onsite potentials and bilinear pair couplings. An internal edge
contains
\begin{equation}
 \frac\lambda4(x-y)^2(x^2+y^2)
 =\frac\lambda4(x^4+y^4-2x^3y+2x^2y^2-2xy^3).
 \label{eq:nonbilinear-scalar-obstruction}
\end{equation}
Its mixed derivative is $\lambda(-3x^2/2+2xy-3y^2/2)$, not a
constant bilinear coupling. This rules out that direct rewriting, not all
possible comparison theorems or reductions. Setting $\lambda=0$ instead
gives eight decoupled scalar crystals and changes the model; restricting
every $q_y$ to the line spanned by $u$ also changes the finite-volume law.

The scalar asymmetric-potential result \cite[Theorem 1.4]{KK-asym}
allows a transition at a nonzero source. KKK Theorems 3.22 and 3.25 are
also scalar results. They cannot supply the present vector conclusion by
citation alone. The candidate contribution is consequently the exact
positive-$\lambda$ collective lower bound and its explicit loop/DLR
composition, not the standard inequalities or threshold shape. Failure of
these particular imports does not establish novelty or publication value;
a broader comparison theorem and specialist assessment remain to be checked.

\section{Reproducibility and adversarial boundary}
\label{sec:repro}

The paper is accompanied by the canonical R-497 source manifest and the
following deterministic finite diagnostics:
\begin{center}
\begin{tabular}{p{0.54\textwidth}p{0.25\textwidth}}
\toprule
Command & Registered result\\
\midrule
\texttt{q3lock\_absolute\_partition\_audit.py} & EXP-001588, exact determinant\\
\texttt{q3lock\_thermodynamic\_pressure\_audit.py} & EXP-001589, seam/tiling\\
\texttt{q3lock\_dlr\_tangent\_content\_audit.py} & EXP-001591, 70/70\\
\texttt{q3lock\_fkg\_content\_audit.py} & EXP-001593, 504/504\\
\texttt{q3lock\_reflection\_infrared\_content\_audit.py} & EXP-001595, finite graph/source checks\\
\texttt{q3lock\_collective\_falk\_bruch\_content\_audit.py} & EXP-001598, 40/40\\
\texttt{q3lock\_strict\_cusp\_tangent\_content\_audit.py} & EXP-001598, 42/42\\
\bottomrule
\end{tabular}
\end{center}
The exact paths, run hashes, and expected outputs are in
\texttt{verification/README.md}.  Each diagnostic was replayed twice with
byte-identical output, but the diagnostics test finite algebra and provenance;
they do not replace a mathematical proof or an external referee.

The adversarial checks that must be signed are:
\begin{enumerate}[leftmargin=2.3em]
\item the corrected seam coefficient $288$ and all source/$\beta$ factors;
\item form-domain membership of the bounded coordinate and the order of
      finite spectral cutoff, coordinate cutoff, and spatial limit;
\item the typed FSS Poisson norm, including the absence of an extra factor of
      eight and the exclusion of the zero mode;
\item transfer of FKG association and fourth-moment products through the
      actual weighted loop limit;
\item the squared-tail estimate used at the nondifferentiable pressure point;
\item compact-uniform source continuity and the clipped expectation passage
      from $h_j\downarrow0$ to a genuine DLR state;
\item the exact sufficient regime, with no inference at equality or outside
      it.
\end{enumerate}

\section{Readiness gates and conclusion}
\label{sec:readiness}

The present source is organized for review, not declared ready for submission.
The release conditions are:
\begin{enumerate}[label=\textup{(G\arabic*)},leftmargin=2.3em]
\item complete remaining source-hypothesis/content review and any resulting repairs, followed by independent
      line-by-line mathematical acceptance of all seven blocks;
\item exact literature hypothesis crosswalk and specialist novelty opinion;
\item bounded claim/result-lineage decision consistent with R-497's T0 scope;
\item final notation, bibliography, nonclaim, and convention freeze;
\item fresh source-hash capture and primary/independent/integrated clean
      snapshot replay with byte-stable outputs;
\item release check on the frozen tree and operator confirmation of commit and
      backup;
\item only then, compilation and visual inspection of the final PDF.
\end{enumerate}
No gate above is silently discharged by the local finite diagnostics.  Until
they are discharged, the scientifically correct description is
\textit{conditional internal manuscript content, independently reviewable,
PDF deferred}.

\begin{thebibliography}{99}
\bibitem{Simon-FK}
B.~Simon,
\emph{A Feynman--Kac Formula for Unbounded Semigroups},
arXiv:math-ph/9907022v1 (1999),
\url{https://arxiv.org/pdf/math-ph/9907022v1}.

\bibitem{KP}
Y.~Kozitsky and T.~Pasurek,
\emph{Euclidean Gibbs Measures of Interacting Quantum Anharmonic Oscillators},
arXiv:math-ph/0609045v1 (2006),
\url{https://arxiv.org/pdf/math-ph/0609045v1}.

\bibitem{FSS}
J.~Fr\"ohlich, B.~Simon, and T.~Spencer,
Infrared bounds, phase transitions and continuous symmetry breaking,
\emph{Commun. Math. Phys.} \textbf{50} (1976), 79--95,
\url{https://math.caltech.edu/SimonPapers/65.pdf}.

\bibitem{FKG}
C.~M.~Fortuin, P.~W.~Kasteleyn, and J.~Ginibre,
Correlation inequalities on some partially ordered sets,
\emph{Commun. Math. Phys.} \textbf{22} (1971), 89--103,
\url{https://math.bme.hu/~balint/oktatas/perkolacio/percolation_papers/fortuin_kasteleyn_ginibre.pdf}.

\bibitem{KKK}
A.~Kargol, Y.~Kondratiev, and Y.~Kozitsky,
\emph{Phase Transitions and Quantum Stabilization in Quantum Anharmonic Crystals},
arXiv:0710.2303v1 (2007),
\url{https://arxiv.org/pdf/0710.2303v1}.

\bibitem{KK-asym}
A.~Kargol and Y.~Kozitsky,
A phase transition in a quantum crystal with asymmetric potentials,
arXiv:math-ph/0611017 (2006),
\url{https://web.ma.utexas.edu/mp_arc/c/06/06-321.pdf}.

\end{thebibliography}

\end{document}
